% =============================================================================
%  MEMORIA TFG — FurniGest
%  Aarón Salinero Valencia — Grado en Ingeniería Informática (URJC)
%  Tutor: Nicolás Rodríguez Uribe — Curso 2024-2025
% =============================================================================
\documentclass[12pt,twoside,titlepage]{report}

%%%%%%%%%%%%%%%%%%%%%%% Paquetes %%%%%%%%%%%%%%%%%%%%%%%
\usepackage[a4paper,bindingoffset=3mm,bottom=35mm]{geometry}
\usepackage[colorlinks=true,pdftex]{hyperref}
\usepackage[pdftex]{graphicx}
\usepackage[spanish]{babel}
\usepackage[utf8]{inputenc}
\usepackage{amsmath,amssymb}
\usepackage{color}
\usepackage{afterpage}
\usepackage{paralist}
\usepackage{array}
\usepackage{enumerate}
\usepackage{enumitem}
\usepackage{float}
\usepackage{setspace}
\usepackage{listings}
\usepackage{algorithm}
\usepackage{algorithmic}
\usepackage{fancyhdr}
\usepackage{rotating}
\usepackage{multirow}
\usepackage{booktabs}
\usepackage{eurosym}
\usepackage{tabularx}
\usepackage{longtable}
\usepackage{quotchap}
\usepackage{tikz}
\usetikzlibrary{shapes.geometric,arrows.meta,positioning,fit,calc,backgrounds,automata,babel}
\usepackage{subcaption}
\usepackage{pgfplots}
\pgfplotsset{compat=1.18}
%%%%%%%%%%%%%%%%%%%%%%%%%%%%%%%%%%%%%%%%%%%%%%%%%%%%%%%%

%%%%%%%%%%%%%%%%%%%%%%% Definiciones básicas %%%%%%%%%%%%%%%%%%%%%%%
\newcommand{\nombreautor}{Aarón Salinero Valencia}
\newcommand{\nombretutor}{Nicolás Rodríguez Uribe}
\newcommand{\titulotrabajo}{FurniGest: aplicación web de gestión integral para tienda de muebles}
\newcommand{\escuela}{Escuela Técnica Superior\\de Ingeniería Informática}
\newcommand{\escuelalargo}{Escuela Técnica Superior de Ingeniería Informática}
\newcommand{\universidad}{Universidad Rey Juan Carlos}
\newcommand{\fecha}{Junio 2025}
\newcommand{\grado}{Grado en Ingeniería Informática}
\newcommand{\curso}{Curso 2024-2025}
\newcommand{\logoUniversidad}{logoURJC.pdf}
%%%%%%%%%%%%%%%%%%%%%%%%%%%%%%%%%%%%%%%%%%%%%%%%%%%%%%%%%%%%%%%%%%%%

%%%%%%%%%%%%%%%%%%%%%%%%% Otras definiciones %%%%%%%%%%%%%%%%%%%%%%%%%%
\definecolor{BlueLink}{rgb}{0.165,0.322,0.745}
\definecolor{PinkLink}{rgb}{0.8,0.22,0.5}
\definecolor{gray}{rgb}{0.6,0.6,0.6}

\hypersetup{hidelinks,pageanchor=true,colorlinks,citecolor=PinkLink,urlcolor=black,linkcolor=BlueLink}

\newcommand\blankpage{%
    \newpage
    \null
    \thispagestyle{empty}%
    \newpage}

\addto{\captionsspanish}{\renewcommand{\bibname}{Bibliografía}}
\addto\captionsspanish{
\def\tablename{Tabla}
\def\listtablename{\'{I}ndice de tablas}
}

\floatname{algorithm}{Algoritmo}
\newfloat{algorithm}{t}{lop}
%%%%%%%%%%%%%%%%%%%%%%%%%%%%%%%%%%%%%%%%%%%%%%%%%%%%%%%%%%%%%%%%%%%%

%%%%%%%%%%%%%%%%%%%%%%% Estilo de código %%%%%%%%%%%%%%%%%%%%%%%
\definecolor{bg}{rgb}{0.95,0.95,0.95}
\definecolor{mydeepteal}{rgb}{0.16,0.22,0.23}
\definecolor{myteal}{rgb}{0.31,0.44,0.46}
\definecolor{mymediumteal}{rgb}{0.41,0.58,0.60}

\DeclareFixedFont{\ttb}{T1}{txtt}{bx}{n}{12}
\DeclareFixedFont{\ttm}{T1}{txtt}{m}{n}{12}
\newcommand*{\FormatDigit}[1]{\textcolor{black}{#1}}

\newcommand\mypythonstyle{\lstset{
language=Python,
basicstyle=\ttfamily\small,
otherkeywords={self},
keywordstyle=\bfseries\ttfamily\color{myteal},
commentstyle=\itshape\color{myteal},
stringstyle=\color{mydeepteal},
emph={MyClass,__init__},
emphstyle=\ttb\color{mydeepteal},
showstringspaces=false,
backgroundcolor=\color{bg},
rulecolor=\color{bg},
breaklines=true,
numbers=left,
numbersep=5pt,
numberstyle=\tiny,
tabsize=4,
xleftmargin=1em,
frame=single,
framesep=3pt,
framextopmargin=0pt,
framexbottommargin=0pt,
framexleftmargin=0pt,
framexrightmargin=0pt,
fontadjust=true,
basewidth=0.55em,
upquote=true,
}}

\lstnewenvironment{mypython}[1][]
{\mypythonstyle\lstset{#1}}
{}

\lstdefinestyle{myjavascript}{
language=Java,
basicstyle=\ttfamily\small,
keywordstyle=\bfseries\ttfamily\color{myteal},
commentstyle=\itshape\color{myteal},
stringstyle=\color{mydeepteal},
showstringspaces=false,
backgroundcolor=\color{bg},
rulecolor=\color{bg},
breaklines=true,
numbers=left,
numbersep=5pt,
numberstyle=\tiny,
tabsize=2,
xleftmargin=1em,
frame=single,
framesep=3pt,
fontadjust=true,
basewidth=0.55em,
upquote=true,
morekeywords={const,let,async,await,export,import,from,default,=>},
}
%%%%%%%%%%%%%%%%%%%%%%%%%%%%%%%%%%%%%%%%%%%%%%%%%%%%%%%%%%%%%%%%%%%%%%

\newcommand{\transpuesta}{\mbox{\tiny $\mathsf{T}$}}

% =============================================================================
%                           INICIO DEL DOCUMENTO
% =============================================================================
\begin{document}
\pagestyle{plain}

%%%%%%%%%%%%%%%%%%%%%%%%%%%%%%%%%%%% Portada %%%%%%%%%%%%%%%%%%%%%%%%%%%%%%%%%%
\begin{titlepage}
\selectlanguage{spanish}

\begin{center}
    \includegraphics[scale=0.7]{\logoUniversidad}
\end{center}

\bigskip

\begin{center}
\begin{LARGE}
\escuela \\
\end{LARGE}
\end{center}

\bigskip\bigskip

\begin{center}
\begin{large}
\textbf{\grado}\\
\end{large}
\end{center}

\begin{center}
\begin{large}
\textbf{\curso}\\
\end{large}
\end{center}

\bigskip

\textbf{\begin{center}
\begin{large}
\textbf{Trabajo Fin de Grado}
\end{large}
\end{center}}

\bigskip\bigskip\bigskip

\begin{center}
\textbf{\begin{large}
\MakeUppercase{\titulotrabajo}\\
\end{large}}
\end{center}

\vspace{\fill}
\begin{center}
\textbf{Autor: \nombreautor}\\ \smallskip
\textbf{Tutor: \nombretutor}\\
\bigskip
\end{center}
\end{titlepage}

\hypersetup{pageanchor=true}
\normalsize
\afterpage{\blankpage}

% Estilo de párrafo
\setlength{\parskip}{0.75em}
\renewcommand{\baselinestretch}{1.25}
\spacing{1}

\pagenumbering{Roman}
\setcounter{page}{2}

% =============================================================================
%                          AGRADECIMIENTOS
% =============================================================================
\chapter*{Agradecimientos}

A mi familia, por su apoyo incondicional durante toda la carrera. A mi tutor, Nicolás Rodríguez Uribe, por su orientación y confianza en este proyecto. A los compañeros que han contribuido con ideas, pruebas y feedback durante el desarrollo de FurniGest.

\afterpage{\blankpage}

% =============================================================================
%                              RESUMEN
% =============================================================================
\chapter*{Resumen}

Este Trabajo Fin de Grado presenta el diseño, implementación y despliegue de \textbf{FurniGest}, un sistema de información de gestión (MIS) para una tienda de muebles. La aplicación centraliza la gestión de clientes, productos, proveedores, albaranes de entrega, movimientos financieros, logística de transporte y analíticas de negocio en una plataforma web accesible desde cualquier navegador.

El \emph{backend} se ha desarrollado con \textbf{FastAPI} (Python 3.12), \textbf{SQLAlchemy 2} como ORM y \textbf{PostgreSQL 16} como sistema gestor de bases de datos. El \emph{frontend} es una aplicación de página única (SPA) construida con \textbf{React 19}, \textbf{Vite 6} y \textbf{TailwindCSS 4}. La comunicación entre ambas capas se realiza mediante una API REST documentada con OpenAPI.

Entre las funcionalidades avanzadas destacan: un motor de analíticas con segmentación RFM de clientes y análisis de cesta (Market Basket Analysis); un asistente de inteligencia artificial basado en Llama-3 a través de Groq; predicciones de ingresos mediante suavizado exponencial de Holt; integración de pagos con Stripe Checkout; generación automática de PDF con ReportLab; envío de correos electrónicos con adjuntos; e internacionalización completa (ES/EN).

Se ha aplicado una estrategia rigurosa de ingeniería del software: más de 540 pruebas automatizadas (unitarias, de integración y end-to-end), análisis estático con Ruff y ESLint, \emph{quality gate} con SonarCloud, un pipeline de integración y despliegue continuos con GitHub Actions, y contenerización con Docker. El sistema está desplegado en producción en Railway (backend) y Vercel (frontend).

\mbox{} \bigskip

\noindent \textbf{Palabras clave}:
\begin{compactitem}
    \item Sistema de información de gestión (MIS)
    \item FastAPI, React, PostgreSQL
    \item Integración continua / despliegue continuo (CI/CD)
    \item Pruebas automatizadas
    \item Inteligencia artificial (LLM)
    \item Análisis de datos (RFM, Market Basket)
\end{compactitem}

\afterpage{\blankpage}

% =============================================================================
%                              ÍNDICES
% =============================================================================
\setlength{\parskip}{1pt}
\renewcommand{\baselinestretch}{1}
\renewcommand{\contentsname}{Índice de contenidos}

\tableofcontents
\afterpage{\blankpage}

\listoftables
\afterpage{\blankpage}
\addcontentsline{toc}{chapter}{\noindent \listtablename}

\listoffigures
\afterpage{\blankpage}
\addcontentsline{toc}{chapter}{\listfigurename}

\renewcommand\lstlistlistingname{Códigos}
\renewcommand\lstlistingname{Código}
\renewcommand\lstlistlistingname{Índice de códigos}
\lstlistoflistings
\afterpage{\blankpage}
\addcontentsline{toc}{chapter}{\lstlistlistingname}

% =============================================================================
%                     CABECERAS Y PIES DE PÁGINA
% =============================================================================
\pagestyle{fancy}
\renewcommand{\chaptermark}[1]{\markboth{Capítulo \thechapter.\ #1}{}}
\fancyhf{}
\fancyhead[LO]{\leftmark}
\fancyhead[RE]{\nouppercase\rightmark}
\fancyfoot[C]{\thepage}

% #############################################################################
%
%                         CAPÍTULO 1: INTRODUCCIÓN
%
% #############################################################################
\chapter{Introducción}
\pagestyle{fancy}
\setlength{\parskip}{0.75em}
\renewcommand{\baselinestretch}{1.25}
\spacing{1}
\pagenumbering{arabic}
\setcounter{page}{1}

\section{Contexto y motivación}

El sector del comercio minorista de muebles en España está formado en gran parte por pequeñas y medianas empresas que gestionan sus operaciones diarias con herramientas genéricas (hojas de cálculo, documentos en papel o software de facturación básico). Esta forma de trabajo presenta problemas recurrentes: pérdida de trazabilidad en los pedidos, duplicación de datos de clientes, imposibilidad de obtener métricas de negocio en tiempo real y comunicación manual con los clientes para confirmar entregas y cobros.

Un \textbf{Sistema de Información de Gestión} (Management Information System, MIS) \cite{pressman} permite capturar, almacenar, procesar y distribuir datos de negocio para apoyar la toma de decisiones. Cuando un MIS se adapta específicamente al dominio de una empresa ---clientes, pedidos, logística, finanzas--- recibe habitualmente el nombre de \textbf{ERP} (Enterprise Resource Planning).

Las soluciones ERP existentes en el mercado (Odoo, SAP Business One, Holded) son potentes pero presentan una barrera significativa para pequeños negocios: complejidad de configuración, costes de licencia elevados y módulos genéricos que no se ajustan al flujo de trabajo específico de una tienda de muebles, donde conceptos como el \emph{albarán de entrega}, la \emph{fianza}, las \emph{rutas de reparto por camión} o la \emph{liquidación de transporte} son centrales.

\textbf{FurniGest} nace como respuesta a esta necesidad: una aplicación web de gestión integral, desarrollada a medida, que cubre desde la creación de pedidos hasta el cobro final, pasando por logística, analíticas avanzadas, inteligencia artificial y personalización de la interfaz. El proyecto se ha concebido desde el inicio con vocación de \textbf{producción real}, no como un prototipo académico: está desplegado en infraestructura cloud, cuenta con integración y despliegue continuos, y cumple estándares de calidad verificados por SonarCloud.

\section{Alcance del proyecto}

FurniGest abarca las siguientes áreas funcionales:

\begin{itemize}[nosep]
  \item Gestión de clientes, productos y proveedores (CRUD completo).
  \item Ciclo de vida del albarán de entrega con máquina de cinco estados.
  \item Registro automático de movimientos financieros (ingresos y egresos).
  \item Módulo de logística: asignación a camiones, manifestos de carga, liquidación de rutas.
  \item Motor de analíticas: segmentación RFM, análisis de cesta, predicción con suavizado exponencial de Holt.
  \item Asistente de inteligencia artificial (Groq / Llama-3).
  \item Cobro con Stripe Checkout.
  \item Generación de PDF (albarán, factura de ruta, informe de tendencias).
  \item Notificación por correo electrónico con PDF adjunto.
  \item Personalización visual (temas, modo oscuro, logotipo de tienda).
  \item Internacionalización ES/EN y accesibilidad WCAG 2.1 AA.
  \item Gestión de incidencias en entregas.
\end{itemize}

Quedan \textbf{fuera del alcance}: contabilidad fiscal y tributaria, gestión de recursos humanos, soporte multi-tienda, aplicación móvil nativa y gestión de inventario/stock en tiempo real.

\section{Estructura del documento}

El resto de esta memoria se organiza de la siguiente forma:

\begin{description}[nosep]
  \item[Capítulo 2 --- Objetivos y planificación.] Define el problema, analiza alternativas, establece los objetivos del proyecto y describe la metodología de desarrollo y la planificación temporal.
  \item[Capítulo 3 --- Análisis y diseño.] Recoge la especificación de requisitos funcionales y no funcionales, la arquitectura del sistema, el modelo de datos, el diseño de la API REST y las decisiones tecnológicas.
  \item[Capítulo 4 --- Implementación.] Detalla la implementación del backend, los flujos de negocio, los algoritmos avanzados, la integración con LLM y el frontend.
  \item[Capítulo 5 --- Pruebas, calidad y despliegue.] Describe la estrategia de pruebas, el análisis estático, la integración continua con GitHub Actions, la contenerización con Docker y el despliegue en producción.
  \item[Capítulo 6 --- Resultados.] Presenta la evaluación de usabilidad, las métricas de calidad del software y las capturas de pantalla del sistema.
  \item[Capítulo 7 --- Conclusiones y trabajos futuros.] Sintetiza las conclusiones del proyecto y propone líneas de trabajo futuro.
\end{description}


% #############################################################################
%
%                    CAPÍTULO 2: OBJETIVOS Y PLANIFICACIÓN
%
% #############################################################################
\chapter{Objetivos y planificación}

\section{Descripción del problema}

Una tienda de muebles de tamaño medio gestiona diariamente decenas de pedidos, cada uno con múltiples productos, datos de cliente, dirección de entrega, estado de pago (fianza, pago pendiente, pago completo) y asignación a un camión de reparto. El flujo típico implica:

\begin{enumerate}[nosep]
  \item El cliente encarga muebles en tienda; se registra un albarán con una fianza.
  \item Los productos se almacenan hasta que se organiza una ruta de reparto.
  \item Un camión recoge los pedidos; el repartidor entrega y confirma.
  \item Se registra el cobro pendiente y se liquidan los costes de transporte.
\end{enumerate}

Cuando este flujo se gestiona con papel y hojas de cálculo surgen problemas concretos: albaranes extraviados, fianzas no registradas, clientes sin actualizar, imposibilidad de saber qué producto se vende más o qué cliente lleva meses sin comprar, y ausencia total de métricas para la toma de decisiones.

\section{Estudio de alternativas}

Antes de optar por un desarrollo a medida, se evaluaron las soluciones existentes recogidas en la tabla~\ref{tab:alternativas}.

\begin{table}[H]
  \centering
  \caption{Comparativa de soluciones existentes.}
  \label{tab:alternativas}
  \begin{tabularx}{\textwidth}{lllX}
    \toprule
    \textbf{Solución} & \textbf{Tipo} & \textbf{Coste} & \textbf{Limitaciones para el dominio} \\
    \midrule
    Odoo        & ERP genérico & Desde 24,90\,\euro/mes & Complejo de configurar; módulo de logística no adaptado a rutas de muebles. \\
    Holded      & SaaS         & Desde 15\,\euro/mes    & Orientado a facturación; sin gestión de rutas, fianzas ni analíticas avanzadas. \\
    Excel / Sheets & Manual    & Gratuito               & Sin trazabilidad, sin automatización, sin escalabilidad. \\
    \textbf{FurniGest} & \textbf{A medida} & \textbf{Desarrollo propio} & \textbf{Adaptado al dominio; control total del código; despliegue en producción.} \\
    \bottomrule
  \end{tabularx}
\end{table}

La solución a medida permite modelar exactamente el flujo de negocio (fianza $\to$ almacén $\to$ ruta $\to$ entrega), incorporar analíticas específicas del sector y mantener el control total sobre la evolución del producto.

\section{Objetivos del proyecto}

\subsection{Objetivo general}

Diseñar, implementar y desplegar en producción una aplicación web de gestión integral para una tienda de muebles que centralice clientes, pedidos, logística, finanzas y analíticas de negocio.

\subsection{Objetivos específicos}

La tabla~\ref{tab:objetivos} enumera los diez objetivos específicos del proyecto.

\begin{table}[H]
  \centering
  \caption{Objetivos específicos del proyecto.}
  \label{tab:objetivos}
  \begin{tabularx}{\textwidth}{lX}
    \toprule
    \textbf{ID} & \textbf{Descripción} \\
    \midrule
    OBJ-01 & Implementar un CRUD completo de clientes, productos y proveedores con búsqueda y filtrado. \\
    OBJ-02 & Modelar el ciclo de vida del albarán como una máquina de estados con cinco transiciones y efectos secundarios automáticos. \\
    OBJ-03 & Registrar automáticamente movimientos financieros (ingresos y egresos) ligados a eventos de negocio. \\
    OBJ-04 & Desarrollar un módulo de logística que permita asignar albaranes a camiones, visualizar manifiestos de carga y liquidar rutas. \\
    OBJ-05 & Construir un motor de analíticas con segmentación RFM de clientes, análisis de cesta y predicción de ingresos (Holt). \\
    OBJ-06 & Integrar un asistente de inteligencia artificial basado en un LLM (Llama-3 vía Groq) con inyección de métricas en tiempo real. \\
    OBJ-07 & Incorporar cobro con tarjeta a través de Stripe Checkout con registro automático del pago como movimiento. \\
    OBJ-08 & Diseñar y ejecutar una suite de pruebas automatizadas que cubra pruebas unitarias, de integración y end-to-end. \\
    OBJ-09 & Configurar un pipeline de integración y despliegue continuos con GitHub Actions y SonarCloud. \\
    OBJ-10 & Desplegar la aplicación en producción en infraestructura cloud (Railway + Vercel). \\
    \bottomrule
  \end{tabularx}
\end{table}

\section{Metodología de desarrollo}

El desarrollo de FurniGest ha seguido un enfoque \textbf{ágil} \cite{manifiestoAgil} basado en iteraciones cortas con entrega incremental de funcionalidad. No se ha aplicado un marco prescriptivo completo como Scrum \cite{schwaber}, sino una adaptación ligera orientada a un equipo de una persona.

\subsection{Trunk-Based Development}

El modelo de ramificación utilizado es \textbf{Trunk-Based Development} (TBD) \cite{tbd}: \texttt{main} es la única rama de larga duración. Cada funcionalidad se desarrolla en una rama \texttt{feature/*} o \texttt{fix/*} de corta vida y se integra a \texttt{main} mediante una Pull Request. Este modelo se adapta al despliegue continuo: cada \emph{push} ejecuta el linting, y cada Pull Request a \texttt{main} ejecuta el pipeline completo de pruebas, quality gate y despliegue.

\begin{figure}[H]
  \centering
  \begin{tikzpicture}[
    node distance=1.2cm and 2cm,
    commit/.style={circle, draw, fill=BlueLink!20, minimum size=8mm, font=\footnotesize},
    branch/.style={rectangle, rounded corners, draw, fill=gray!20, minimum height=7mm, font=\footnotesize},
    >=Stealth
  ]
    \node[commit] (m1) {m1};
    \node[commit, right=of m1] (m2) {m2};
    \node[commit, right=of m2] (m3) {m3};
    \node[commit, right=of m3] (m4) {m4};
    \node[commit, right=of m4] (m5) {m5};

    \draw[->] (m1) -- (m2);
    \draw[->] (m2) -- (m3);
    \draw[->] (m3) -- (m4);
    \draw[->] (m4) -- (m5);

    \node[commit, fill=green!20, above right=0.8cm and 0.8cm of m1] (f1) {f1};
    \node[commit, fill=green!20, right=of f1] (f2) {f2};

    \draw[->, dashed] (m1) -- (f1);
    \draw[->] (f1) -- (f2);
    \draw[->, dashed] (f2) -- (m3);

    \node[commit, fill=orange!20, above right=0.8cm and 0.8cm of m3] (g1) {g1};
    \draw[->, dashed] (m3) -- (g1);
    \draw[->, dashed] (g1) -- (m5);

    \node[branch, below=0.4cm of m3] {\texttt{main}};
    \node[branch, fill=green!10, above=0.3cm of f1] {\texttt{feature/analytics}};
    \node[branch, fill=orange!10, above=0.3cm of g1] {\texttt{fix/email}};
  \end{tikzpicture}
  \caption{Trunk-Based Development: rama \texttt{main} con ramas de corta vida.}
  \label{fig:tbd}
\end{figure}

\subsection{Conventional Commits}

Todos los mensajes de commit siguen la especificación \textbf{Conventional Commits} \cite{conventionalcommits}: \texttt{<tipo>[ámbito]: <descripción>}. Los tipos utilizados son: \texttt{feat} (nueva funcionalidad), \texttt{fix} (corrección), \texttt{refactor} (reestructuración sin cambio de comportamiento), \texttt{test} (pruebas), \texttt{docs} (documentación) y \texttt{chore} (mantenimiento).

\section{Planificación temporal}

La tabla~\ref{tab:planificacion} resume las fases del proyecto y su duración aproximada.

\begin{table}[H]
  \centering
  \caption{Planificación temporal del proyecto.}
  \label{tab:planificacion}
  \begin{tabular}{lc}
    \toprule
    \textbf{Fase} & \textbf{Duración} \\
    \midrule
    Definición de requisitos y diseño de arquitectura & 3 semanas \\
    Implementación del backend (funcionalidades base) & 4 semanas \\
    Implementación del frontend (funcionalidades base) & 3 semanas \\
    Funcionalidades avanzadas (analíticas, IA, Stripe, i18n) & 3 semanas \\
    Pruebas automatizadas y configuración CI/CD & 2 semanas \\
    Despliegue en producción y documentación & 2 semanas \\
    \midrule
    \textbf{Total} & \textbf{17 semanas} \\
    \bottomrule
  \end{tabular}
\end{table}

\blankpage


% #############################################################################
%
%                      CAPÍTULO 3: ANÁLISIS Y DISEÑO
%
% #############################################################################
\chapter{Análisis y diseño}

\section{Especificación de requisitos}

La especificación de requisitos sigue las recomendaciones del estándar IEEE~830 \cite{ieee830}.

\subsection{Requisitos funcionales}

La tabla~\ref{tab:rf} recoge los dieciocho requisitos funcionales del sistema.

\begin{longtable}{lp{11cm}}
  \caption{Requisitos funcionales.} \label{tab:rf} \\
  \toprule
  \textbf{ID} & \textbf{Descripción} \\
  \midrule
  \endfirsthead
  \toprule
  \textbf{ID} & \textbf{Descripción} \\
  \midrule
  \endhead
  RF-01 & Gestión de clientes: crear, leer, actualizar y eliminar clientes. Búsqueda por DNI y correo electrónico con \emph{upsert} automático. \\
  RF-02 & Gestión de productos y proveedores: CRUD de ambas entidades con relación uno a muchos (un proveedor tiene muchos productos). \\
  RF-03 & Ciclo de vida del albarán: máquina de cinco estados (FIANZA $\to$ ALMACEN $\to$ RUTA $\to$ ENTREGADO $\to$ INCIDENCIA) con transiciones controladas. \\
  RF-04 & Libro de movimientos financieros: registro y listado de ingresos y egresos, con generación automática al crear albaranes y al entregar pedidos. \\
  RF-05 & Dashboard: pantalla resumen con indicadores clave de rendimiento (ingresos totales, pedidos abiertos, clientes activos) y gráficos interactivos. \\
  RF-06 & Notificación por email: envío automático de correo electrónico con PDF adjunto del albarán al crear un pedido. \\
  RF-07 & Motor de analíticas: cálculo de ventas diarias, semanales, comparativa entre períodos y ranking de productos por ingresos. \\
  RF-08 & Segmentación RFM: clasificación automática de clientes en segmentos (VIP, En crecimiento, En riesgo, Ocasional) según recencia, frecuencia y valor monetario. \\
  RF-09 & Análisis de cesta: descubrimiento de pares de productos comprados frecuentemente juntos, con cálculo de soporte, confianza y lift. \\
  RF-10 & Asistente IA: interfaz conversacional con un LLM que recibe métricas de negocio en tiempo real y genera respuestas y gráficos. \\
  RF-11 & Módulo de transporte: asignación de albaranes a camiones, visualización de manifiestos de carga, liquidación de rutas (7\,\% de coste) y descarga de factura PDF. \\
  RF-12 & Cobro con Stripe: generación de sesión de Stripe Checkout, redirección al pago, confirmación y registro automático del ingreso. \\
  RF-13 & Generación de PDF: exportación de albarán, factura de ruta e informe de tendencias en formato PDF. \\
  RF-14 & Personalización visual: seis paletas de color, modo oscuro/claro, logotipo y nombre de tienda configurables. \\
  RF-15 & Resumen semanal: envío periódico automático por email de un resumen de negocio generado por IA, con frecuencia e intervalo configurables. \\
  RF-16 & Perfil y configuración: actualización de contraseña y nombre de usuario; almacén clave-valor para parámetros de la tienda. \\
  RF-17 & Internacionalización: interfaz completa en español e inglés, con cambio de idioma en tiempo de ejecución. \\
  RF-18 & Gestión de incidencias: registro de problemas de entrega vinculados a un albarán; al eliminar la incidencia, el albarán revierte a ENTREGADO. \\
  RF-19 & Integración bancaria (CaixaBank PSD2): enlace OAuth con Open Banking, consulta de cuentas y estado de conexión, como base para la futura conciliación automática de pagos. \\
  \bottomrule
\end{longtable}

\subsection{Requisitos no funcionales}

La tabla~\ref{tab:rnf} especifica los requisitos no funcionales del sistema.

\begin{table}[H]
  \centering
  \caption{Requisitos no funcionales.}
  \label{tab:rnf}
  \begin{tabularx}{\textwidth}{lX}
    \toprule
    \textbf{ID} & \textbf{Descripción} \\
    \midrule
    RNF-01 & \textbf{Rendimiento:} tiempo de respuesta de la API inferior a 200\,ms en operaciones CRUD. \\
    RNF-02 & \textbf{Disponibilidad:} despliegue en Railway con reinicio automático ante fallos y \emph{healthcheck}. \\
    RNF-03 & \textbf{Seguridad:} autenticación JWT \cite{rfc7519}, hashing de contraseñas con bcrypt, CORS explícito, TLS en producción. Alineación con OWASP Top~10 \cite{owasp2021}. \\
    RNF-04 & \textbf{Mantenibilidad:} linting con Ruff y ESLint; cobertura de pruebas superior al 80\,\%; \emph{quality gate} SonarCloud. \\
    RNF-05 & \textbf{Portabilidad:} contenerización con Docker Compose (tres servicios); ejecutable en cualquier máquina con Docker. \\
    RNF-06 & \textbf{Usabilidad:} diseño centrado en el usuario; cumplimiento WCAG 2.1 nivel AA; evaluación heurística de Nielsen \cite{nielsen}. \\
    RNF-07 & \textbf{Escalabilidad:} arquitectura \emph{stateless} en el backend; base de datos relacional con índices en claves foráneas y campos de filtrado. \\
    RNF-08 & \textbf{Internacionalización:} interfaz en español e inglés con detección automática de idioma. \\
    \bottomrule
  \end{tabularx}
\end{table}

\subsection{Perfiles de usuario}

Siguiendo la metodología de \textbf{Diseño Centrado en el Usuario} (UCD) \cite{iso9241}, se han identificado tres perfiles de usuario (personas) recogidos en la tabla~\ref{tab:personas}.

\begin{table}[H]
  \centering
  \caption{Perfiles de usuario (personas UCD).}
  \label{tab:personas}
  \begin{tabularx}{\textwidth}{llX}
    \toprule
    \textbf{Persona} & \textbf{Rol} & \textbf{Necesidades clave} \\
    \midrule
    Gerente de tienda & Visión estratégica & Dashboard, analíticas, segmentación RFM, asistente IA, informes PDF. \\
    Vendedor & Operaciones diarias & Crear pedidos, buscar clientes y productos, avanzar estados del albarán. \\
    Repartidor & Logística & Manifiesto de carga por camión, cambio de estado en ruta, liquidación de ruta. \\
    \bottomrule
  \end{tabularx}
\end{table}


\section{Arquitectura del sistema}

\subsection{Arquitectura cliente-servidor}

FurniGest sigue una arquitectura \textbf{cliente-servidor} de tres capas \cite{tanenbaum}: un cliente ligero (navegador web con SPA en React), un servidor de aplicaciones (API REST con FastAPI) y un servidor de base de datos (PostgreSQL). La comunicación entre cliente y servidor se realiza mediante el protocolo \textbf{HTTP/1.1} \cite{kurose} sobre TLS en producción, siguiendo el estilo arquitectónico \textbf{REST} (Representational State Transfer) definido por Fielding \cite{fieldingREST}.

\begin{figure}[H]
  \centering
  \begin{tikzpicture}[
    box/.style={rectangle, draw, rounded corners, minimum width=3.5cm, minimum height=1.5cm, align=center, font=\small},
    >=Stealth
  ]
    \node[box, fill=blue!10]  (fe)  {React 19 SPA\\(Vite 6 + TailwindCSS 4)};
    \node[box, fill=green!10, right=2.5cm of fe] (be)  {FastAPI 0.115\\(Python 3.12)};
    \node[box, fill=orange!10, right=2.5cm of be] (db)  {PostgreSQL 16};

    \draw[<->, thick] (fe) -- node[above, font=\footnotesize] {HTTP/JSON} node[below, font=\footnotesize] {REST API} (be);
    \draw[<->, thick] (be) -- node[above, font=\footnotesize] {SQL} node[below, font=\footnotesize] {SQLAlchemy 2} (db);

    \node[below=0.2cm of fe, font=\footnotesize\itshape] {Cliente};
    \node[below=0.2cm of be, font=\footnotesize\itshape] {Servidor de aplicaciones};
    \node[below=0.2cm of db, font=\footnotesize\itshape] {Base de datos};
  \end{tikzpicture}
  \caption{Arquitectura de tres capas de FurniGest.}
  \label{fig:arquitectura_3capas}
\end{figure}

\subsection{Arquitectura interna del backend}

El backend implementa una \textbf{arquitectura en capas} que separa responsabilidades:

\begin{figure}[H]
  \centering
  \begin{tikzpicture}[
    layer/.style={rectangle, draw, minimum width=10cm, minimum height=1cm, align=center, font=\small},
    >=Stealth
  ]
    \node[layer, fill=blue!10]   (api)    {Capa API (routers FastAPI)};
    \node[layer, fill=green!10, below=0.3cm of api]  (svc)    {Capa de servicios (lógica de negocio)};
    \node[layer, fill=yellow!10, below=0.3cm of svc] (entity) {Capa de entidades (ORM SQLAlchemy + schemas Pydantic)};
    \node[layer, fill=orange!10, below=0.3cm of entity] (db)    {Capa de persistencia (PostgreSQL)};

    \draw[->] (api) -- (svc);
    \draw[->] (svc) -- (entity);
    \draw[->] (entity) -- (db);
  \end{tikzpicture}
  \caption{Arquitectura en capas del backend.}
  \label{fig:capas_backend}
\end{figure}

En esta arquitectura se aplican varios \textbf{patrones de diseño} \cite{fowlerRefactoring}, recogidos en la tabla~\ref{tab:patrones}.

\begin{table}[H]
  \centering
  \caption{Patrones de diseño aplicados en el proyecto.}
  \label{tab:patrones}
  \begin{tabularx}{\textwidth}{lX}
    \toprule
    \textbf{Patrón} & \textbf{Aplicación en FurniGest} \\
    \midrule
    Inyección de Dependencias & \texttt{Depends(get\_db)} y \texttt{Depends(get\_current\_user)} inyectan la sesión de BD y el usuario autenticado en cada \emph{endpoint}. \\
    Repository & La sesión de SQLAlchemy encapsula el acceso a datos; los servicios no ejecutan SQL directamente. \\
    DTO (Data Transfer Object) & Pydantic define esquemas separados para creación (\texttt{CustomerCreate}) y respuesta (\texttt{Customer}), desacoplando la representación HTTP del modelo ORM. \\
    Strategy & El módulo de email selecciona en tiempo de ejecución entre \texttt{\_send\_via\_smtp()} y \texttt{\_send\_via\_resend()} según la variable \texttt{RESEND\_API\_KEY}. \\
    Facade & El módulo \texttt{analytics} agrupa RFM, cesta y predicción tras una fachada única (\texttt{/api/analytics/summary}). \\
    \bottomrule
  \end{tabularx}
\end{table}

\subsection{Arquitectura del frontend}

El frontend es una SPA (Single Page Application) con \textbf{React Router 7} \cite{reactrouter}, que monta un \emph{shell} persistente (\texttt{App.jsx} con \texttt{Sidebar.jsx}) y conmuta únicamente el área de contenido mediante \texttt{<Outlet>}. El estado global se gestiona con la \textbf{Context API} de React: \texttt{ThemeContext} (modo oscuro y paleta de colores) y \texttt{ConfigContext} (nombre de tienda, logotipo). Los \emph{hooks} personalizados (\texttt{useClientes}, \texttt{useProductos}, etc.) encapsulan las llamadas a la API.


\section{Modelo de datos}

\subsection{Modelo Entidad-Relación}

El modelo de datos de FurniGest consta de once entidades persistidas en PostgreSQL \cite{postgresql}. El diagrama Entidad-Relación \cite{chen1976} de la figura~\ref{fig:er} muestra las entidades principales y sus relaciones.

\begin{figure}[H]
  \centering
  \begin{tikzpicture}[
    entity/.style={rectangle, draw, fill=blue!8, minimum width=2.8cm, minimum height=0.9cm, font=\small\bfseries, align=center},
    rel/.style={diamond, draw, fill=yellow!15, aspect=2, font=\scriptsize, align=center, inner sep=1pt},
    >=Stealth,
    every edge/.style={draw, -}
  ]
    % Entidades
    \node[entity] (cli) at (0,0) {CustomerDB\\{\scriptsize(clientes)}};
    \node[entity] (alb) at (5,0) {DeliveryNoteDB\\{\scriptsize(albaranes)}};
    \node[entity] (lin) at (5,-2.5) {DeliveryNote\\LineDB\\{\scriptsize(lineas\_albaran)}};
    \node[entity] (pro) at (10,-2.5) {ProductDB\\{\scriptsize(productos)}};
    \node[entity] (sup) at (10,-5) {SupplierDB\\{\scriptsize(proveedores)}};
    \node[entity] (rut) at (10,0) {DeliveryNote\\RouteDB\\{\scriptsize(albaran\_rutas)}};
    \node[entity] (mov) at (0,-2.5) {MovementDB\\{\scriptsize(movimientos)}};
    \node[entity] (inc) at (0,-5) {IncidenciaDB\\{\scriptsize(incidencias)}};
    \node[entity] (str) at (5,-5) {StripeCheckout\\DB};
    \node[entity] (usr) at (0,2) {UserDB\\{\scriptsize(usuarios)}};
    \node[entity] (cfg) at (5,2) {ConfigDB\\{\scriptsize(configuracion)}};

    % Relaciones
    \draw[->] (alb) -- node[above, font=\scriptsize] {N:1} (cli);
    \draw[->] (lin) -- node[left, font=\scriptsize] {N:1} (alb);
    \draw[->] (lin) -- node[above, font=\scriptsize] {N:1} (pro);
    \draw[->] (pro) -- node[left, font=\scriptsize] {N:1} (sup);
    \draw[->] (rut) -- node[above, font=\scriptsize] {1:1} (alb);
    \draw[->, dashed] (inc) -- node[left, font=\scriptsize] {N:1} (alb);
  \end{tikzpicture}
  \caption{Diagrama Entidad-Relación simplificado de FurniGest.}
  \label{fig:er}
\end{figure}

\subsection{Esquema relacional y normalización}

El esquema relacional (modelo de Codd \cite{codd1970}) está en \textbf{Tercera Forma Normal} (3NF):

\begin{itemize}[nosep]
  \item \textbf{1NF:} todos los atributos son atómicos (no hay campos multivaluados).
  \item \textbf{2NF:} cada atributo no clave depende funcionalmente de la clave primaria completa (no hay dependencias parciales).
  \item \textbf{3NF:} no existen dependencias transitivas; por ejemplo, los datos del proveedor no se almacenan en la tabla de productos sino que se referencian mediante \texttt{supplier\_id} (clave foránea).
\end{itemize}

Un caso relevante de diseño es el campo \texttt{unit\_price} en \texttt{lineas\_albaran}: almacena el precio del producto en el momento de la venta, independientemente de cambios posteriores en el catálogo. Esto preserva la integridad del histórico de ventas.

La tabla~\ref{tab:schema_resumen} resume las tablas principales del esquema.

\begin{table}[H]
  \centering
  \caption{Resumen del esquema relacional (tablas principales).}
  \label{tab:schema_resumen}
  \begin{tabular}{llp{6cm}}
    \toprule
    \textbf{Tabla} & \textbf{Clase ORM} & \textbf{Campos clave} \\
    \midrule
    \texttt{clientes}        & \texttt{CustomerDB}         & id, nombre, apellidos, dni, email, teléfono, dirección \\
    \texttt{proveedores}     & \texttt{SupplierDB}         & id, nombre, contacto \\
    \texttt{productos}       & \texttt{ProductDB}          & id, nombre, descripción, precio, proveedor\_id (FK) \\
    \texttt{albaranes}       & \texttt{DeliveryNoteDB}     & id, cliente\_id (FK), fecha, descripción, total, estado \\
    \texttt{lineas\_albaran} & \texttt{DeliveryNoteLineDB} & id, albaran\_id (FK), producto\_id (FK), cantidad, precio\_unitario \\
    \texttt{movimientos}     & \texttt{MovementDB}         & id, fecha, concepto, cantidad, tipo (INGRESO/EGRESO) \\
    \texttt{albaran\_rutas}  & \texttt{DeliveryNoteRouteDB}& id, albaran\_id (FK, unique), camion\_id \\
    \texttt{incidencias}     & \texttt{IncidenciaDB}       & id, albaran\_id (FK), descripción \\
    \texttt{stripe\_checkouts}&\texttt{StripeCheckoutDB}   & id, session\_id (unique), amount, status \\
    \texttt{usuarios}        & \texttt{UserDB}             & id, username (unique), hashed\_password, role \\
    \texttt{configuracion}   & \texttt{ConfigDB}           & clave (PK), valor \\
    \bottomrule
  \end{tabular}
\end{table}

\textbf{Convención de nombrado:} los atributos Python del ORM usan nombres en inglés (\texttt{name}, \texttt{date}, \texttt{amount}), mientras que las columnas de PostgreSQL mantienen sus nombres originales en español (\texttt{nombre}, \texttt{fecha}, \texttt{cantidad}). SQLAlchemy realiza el \emph{alias} automáticamente: \texttt{name = Column('nombre', String)}.

\subsection{PostgreSQL y el teorema CAP}

PostgreSQL es un sistema \textbf{CP} bajo el teorema CAP \cite{brewer2000}: prioriza \textbf{consistencia} y \textbf{tolerancia a particiones} sobre disponibilidad. Para un sistema de gestión interno donde la integridad de datos financieros es crítica (saldos, fianzas, liquidaciones), la consistencia fuerte es la elección adecuada frente a sistemas AP que ofrecen consistencia eventual.


\section{Diseño de la API REST}

La API REST se sirve en \texttt{http://localhost:8000} con prefijo \texttt{/api}. Todos los \emph{endpoints} requieren un token JWT en la cabecera \texttt{Authorization: Bearer <token>}, excepto \texttt{POST /api/auth/login} y \texttt{GET /health}. La documentación interactiva se genera automáticamente con OpenAPI (accesible en \texttt{/docs}).

La tabla~\ref{tab:endpoints} resume los endpoints agrupados por recurso.

\begin{table}[H]
  \centering
  \caption{Resumen de endpoints de la API REST.}
  \label{tab:endpoints}
  \begin{tabular}{lll}
    \toprule
    \textbf{Recurso} & \textbf{Endpoints} & \textbf{Operaciones principales} \\
    \midrule
    Clientes       & 5 & CRUD + búsqueda por DNI/email \\
    Productos      & 6 & CRUD + obtener uno + búsqueda \\
    Proveedores    & 5 & CRUD + obtener uno \\
    Albaranes      & 8 & Crear, listar, obtener, actualizar, cambiar estado, por cliente, editar líneas, PDF \\
    Movimientos    & 5 & CRUD completo + listar con filtro por tipo \\
    Transporte     & 7 & Almacén, en ruta, rutas agrupadas, asignar, quitar, pendiente, liquidar \\
    Analíticas     & 4 & Resumen, comparativa, predicción, exportar PDF \\
    IA             & 2 & Pregunta puntual y conversación multi-turno \\
    Stripe         & 4 & Estado, crear sesión, confirmar pago, listar cobros \\
    Banca (CaixaBank) & 5 & Debug, estado, vincular, callback OAuth, cuentas \\
    Incidencias    & 4 & Listar, obtener, crear, eliminar \\
    Auth           & 3 & Login, obtener perfil, actualizar perfil \\
    Config         & 2 & Obtener y actualizar configuración \\
    Health         & 1 & Sonda de salud \\
    \midrule
    \textbf{Total} & \textbf{61} & \\
    \bottomrule
  \end{tabular}
\end{table}


\section{Decisiones tecnológicas}

Cada tecnología se ha elegido de forma deliberada. La tabla~\ref{tab:stack} recoge el \emph{stack} tecnológico con su justificación; la tabla~\ref{tab:alternativas_tech} las alternativas descartadas.

\begin{table}[H]
  \centering
  \caption{Stack tecnológico y justificación.}
  \label{tab:stack}
  \begin{tabularx}{\textwidth}{llX}
    \toprule
    \textbf{Capa} & \textbf{Tecnología} & \textbf{Justificación} \\
    \midrule
    Backend        & FastAPI 0.115 \cite{fastapi}     & ASGI nativo, documentación OpenAPI automática desde \emph{type hints}, integración con Pydantic. \\
    ORM            & SQLAlchemy 2.x \cite{sqlalchemy}  & Mapeo declarativo, \texttt{relationship()} con cascada, \texttt{selectinload()} para evitar N+1. \\
    Validación     & Pydantic v2 \cite{pydantic}       & \texttt{model\_dump()}, \texttt{from\_attributes=True}, \texttt{Literal[]} para máquina de estados. \\
    Base de datos  & PostgreSQL 16 \cite{postgresql}   & ACID, claves foráneas, JSON nativo, sistema CP. \\
    Frontend       & React 19 \cite{react}             & Hooks, ecosistema (react-chartjs-2, react-router-dom). \\
    Build          & Vite 6 \cite{vite}                & ES modules nativos, HMR < 50\,ms. \\
    CSS            & TailwindCSS 4 \cite{tailwindcss}  & Utility-first, CSS generado en build, cero \emph{runtime}. \\
    Gráficos       & Chart.js 4 \cite{chartjs}         & Renderizado \texttt{<canvas>}, registro selectivo de componentes. \\
    PDF            & ReportLab 4 \cite{reportlab}      & In-process, sin dependencia de navegador. \\
    LLM            & Groq / Llama-3.1 \cite{groq}      & Latencia sub-segundo, API compatible con OpenAI, coste cero en desarrollo. \\
    Pagos          & Stripe Checkout \cite{stripe}     & PCI \emph{offloading}, sesión servidor, idempotencia. \\
    Email          & SMTP + Resend \cite{resend}       & Dual: SMTP local, Resend HTTP en producción (Railway bloquea SMTP). \\
    \bottomrule
  \end{tabularx}
\end{table}

\begin{table}[H]
  \centering
  \caption{Alternativas descartadas con motivo.}
  \label{tab:alternativas_tech}
  \begin{tabularx}{\textwidth}{llX}
    \toprule
    \textbf{Capa} & \textbf{Alternativa} & \textbf{Motivo de descarte} \\
    \midrule
    Backend   & Django REST       & ORM y admin innecesarios para API; overhead excesivo. \\
    Backend   & Flask             & Demasiado cableado manual; sin validación integrada. \\
    ORM       & Tortoise ORM      & Ecosistema más pequeño; menor soporte de \emph{eager loading}. \\
    BD        & SQLite            & Bloqueo de escritura con múltiples conexiones async. \\
    Frontend  & Vue 3             & Ecosistema de gráficos y enrutamiento menos maduro para este caso. \\
    Build     & Webpack / CRA     & CRA sin mantenimiento; Webpack con sobrecarga de configuración. \\
    CSS       & Bootstrap 5       & Clases prediseñadas limitan la personalización por paletas. \\
    LLM       & OpenAI GPT-4o     & Coste por token; modelo 8B suficiente para razonamiento sobre JSON inyectado. \\
    \bottomrule
  \end{tabularx}
\end{table}


% #############################################################################
%
%                      CAPÍTULO 4: IMPLEMENTACIÓN
%
% #############################################################################
\chapter{Implementación}

\section{Estructura del proyecto}

El repositorio se organiza en tres directorios principales: \texttt{backend/} (FastAPI + SQLAlchemy), \texttt{frontend/} (React SPA) y \texttt{test/} (pruebas unitarias, de integración y E2E). Los ficheros de configuración con secretos (\texttt{stripe\_config.py}, \texttt{ia\_settings.py}, \texttt{settings\_email.py}, etc.) están excluidos del repositorio mediante \texttt{.gitignore}.

\section{Backend --- API REST con FastAPI}

\subsection{Punto de entrada y ciclo de vida}

El fichero \texttt{main.py} configura la aplicación FastAPI con un \emph{lifespan} asíncrono que, al arrancar: (1) crea las tablas con \texttt{Base.metadata.create\_all()}, (2) ejecuta el \emph{seed} de datos de demostración (idempotente), y (3) inicia el planificador APScheduler para el resumen semanal. El middleware CORS se configura con los orígenes permitidos (localhost para desarrollo, dominios de Vercel para producción). Cada dominio funcional tiene su propio router registrado con \texttt{app.include\_router()}.

\subsection{Capa de entidades (ORM + schemas)}

Cada entidad tiene un fichero en \texttt{entidades/} con dos elementos: el \textbf{modelo ORM} (SQLAlchemy) y los \textbf{esquemas Pydantic} de validación.

\begin{mypython}[caption={Ejemplo: modelo ORM y schema Pydantic de un cliente.},label={code:cliente}]
# backend/app/entidades/cliente.py
class CustomerDB(Base):
    __tablename__ = "clientes"
    id       = Column(Integer, primary_key=True, index=True)
    name     = Column("nombre", String(100))
    surnames = Column("apellidos", String(150))
    dni      = Column("dni", String(20), index=True)
    email    = Column("email", String(200))
    # ... (direccion, telefonos)

class CustomerCreate(BaseModel):
    name: str
    surnames: str
    dni: str | None = None
    email: str | None = None

class Customer(CustomerCreate):
    id: int
    model_config = ConfigDict(from_attributes=True)
\end{mypython}

La opción \texttt{from\_attributes=True} permite construir el schema Pydantic directamente desde una instancia ORM: \texttt{Customer.model\_validate(customer\_db)}.

\subsection{Capa de servicios}

La lógica de negocio se concentra en \texttt{services/}, separándola de los controladores HTTP. Por ejemplo, \texttt{clientes\_service.upsert\_customer()} busca un cliente por DNI o email antes de crear uno nuevo, evitando duplicados. Los routers de la capa API invocan a los servicios y devuelven la respuesta HTTP.

\subsection{Autenticación y autorización}

La autenticación se basa en \textbf{JSON Web Tokens} (JWT, RFC 7519) \cite{rfc7519}:

\begin{enumerate}[nosep]
  \item El usuario envía \texttt{POST /api/auth/login} con \texttt{username} y \texttt{password}.
  \item El servidor verifica la contraseña contra el hash bcrypt almacenado.
  \item Si es correcta, genera un JWT firmado con HS256 y expiración de 60 minutos.
  \item Todos los \emph{endpoints} protegidos reciben el token en la cabecera \texttt{Authorization: Bearer <token>} y lo validan mediante \texttt{Depends(get\_current\_user)}.
\end{enumerate}

Las contraseñas se almacenan hasheadas con \textbf{bcrypt} (passlib), nunca en texto claro. Esto cubre las recomendaciones A02 (Cryptographic Failures) y A07 (Identification and Authentication Failures) de OWASP Top~10 \cite{owasp2021}.


\section{Flujos de negocio principales}

\subsection{Ciclo de vida del albarán}

El albarán de entrega es la entidad central del sistema. Su ciclo de vida se modela como una \textbf{máquina de estados finita} con cinco estados y transiciones controladas por el backend (figura~\ref{fig:estados_albaran}).

\begin{figure}[H]
  \centering
  \begin{tikzpicture}[
    state/.style={rectangle, rounded corners, draw, minimum width=2.5cm, minimum height=1cm, align=center, font=\small\bfseries},
    >=Stealth,
    every edge/.style={draw, ->},
    note/.style={font=\scriptsize\itshape, text width=3cm, align=center}
  ]
    \node[state, fill=yellow!20]  (fi) {FIANZA};
    \node[state, fill=blue!15, right=2.5cm of fi]   (al) {ALMACEN};
    \node[state, fill=violet!15, right=2.5cm of al]  (ru) {RUTA};
    \node[state, fill=green!15, below=2cm of ru]     (en) {ENTREGADO};
    \node[state, fill=red!15, below=2cm of al]       (inc) {INCIDENCIA};

    \draw[->] (fi) -- node[above, note] {Confirmar pago fianza} (al);
    \draw[->] (al) -- node[above, note] {Asignar a camión} (ru);
    \draw[->] (ru) -- node[right, note] {Confirmar entrega} (en);
    \draw[->] (en) -- node[above, note] {Reportar\\problema} (inc);
    \draw[->, dashed] (inc) -- node[left, note] {Eliminar\\incidencia} (en);

    \node[note, below=0.3cm of fi] {Auto: movimiento\\INGRESO (fianza)};
    \node[note, below=0.3cm of en, text width=3.5cm] {Auto: movimiento INGRESO\\(pendiente = total $-$ fianza)};

  \end{tikzpicture}
  \caption{Máquina de estados del albarán de entrega.}
  \label{fig:estados_albaran}
\end{figure}

Las transiciones generan \textbf{efectos secundarios} automáticos:
\begin{itemize}[nosep]
  \item Al crear el albarán (estado FIANZA): se registra un movimiento de tipo INGRESO por el importe de la fianza (30\,\% por defecto).
  \item Al pasar a ENTREGADO: se registra un segundo movimiento INGRESO por el importe pendiente (total $-$ fianza).
  \item La liquidación de ruta registra un movimiento EGRESO del 7\,\% del total transportado.
\end{itemize}

Adicionalmente, al crear un albarán se encola un \texttt{BackgroundTask} de FastAPI que genera el PDF con ReportLab y envía el correo electrónico al cliente de forma asíncrona, sin bloquear la respuesta HTTP.

\subsection{Flujo de pago con Stripe}

El cobro con tarjeta sigue el flujo estándar de Stripe Checkout \cite{stripe}:

\begin{figure}[H]
  \centering
  \begin{tikzpicture}[
    box/.style={rectangle, draw, rounded corners, minimum width=2.5cm, minimum height=0.9cm, align=center, font=\small},
    >=Stealth
  ]
    \node[box, fill=blue!10]   (fe)     {Frontend\\(BancoPage)};
    \node[box, fill=green!10, right=2cm of fe]  (be)     {Backend\\(FastAPI)};
    \node[box, fill=orange!10, right=2cm of be] (st)     {Stripe\\API};

    \draw[->] (fe) -- node[above, font=\scriptsize] {1. checkout} (be);
    \draw[->] (be) -- node[above, font=\scriptsize] {2. Session.create} (st);
    \draw[<-] (be) -- node[below, font=\scriptsize] {3. session\_id} (st);
    \draw[<-] (fe) -- node[below, font=\scriptsize] {4. redirect URL} (be);

    \node[box, fill=orange!10, below=1.5cm of st] (st2) {Stripe\\Checkout};
    \draw[->] (fe) |- node[right, font=\scriptsize, pos=0.25] {5. redirect} (st2);
    \draw[->] (st2) -| node[right, font=\scriptsize, pos=0.75] {6. success\_url} (fe);

    \draw[->] ([yshift=-3mm]fe.east) -- node[below, font=\scriptsize] {7. confirm} ([yshift=-3mm]be.west);
    \node[font=\scriptsize\itshape, below=3.5cm of be, text width=5cm, align=center] {8. Verificar pago $\to$ registrar movimiento INGRESO};
  \end{tikzpicture}
  \caption{Flujo de pago con Stripe Checkout.}
  \label{fig:stripe_flow}
\end{figure}

La tabla \texttt{stripe\_checkouts} actúa como guarda de idempotencia: si el usuario refresca la página de éxito, el segundo \texttt{confirm} detecta que el \texttt{session\_id} ya existe y no registra un movimiento duplicado.

\subsection{Flujo de email y PDF}

Al crear un albarán, el \texttt{BackgroundTask} ejecuta:
\begin{enumerate}[nosep]
  \item Generación del PDF con ReportLab (\texttt{generate\_delivery\_note\_pdf}).
  \item Renderizado del cuerpo HTML con Jinja2.
  \item Envío según el entorno: \texttt{\_send\_via\_smtp()} (SMTP sobre STARTTLS, puerto 587) o \texttt{\_send\_via\_resend()} (HTTP API sobre puerto 443 cuando \texttt{RESEND\_API\_KEY} está definida).
\end{enumerate}

Este diseño aplica el \textbf{patrón Strategy}: la selección del mecanismo de envío es transparente para el código llamante.


\section{Algoritmos avanzados}

\subsection{Segmentación RFM de clientes}

El modelo RFM \cite{hughes1994,rfm_review} puntúa cada cliente en tres dimensiones extraídas de la tabla \texttt{albaranes}:

\begin{itemize}[nosep]
  \item \textbf{Recency (R):} días desde el último pedido ($R = (\text{hoy} - \max(\text{fecha})).\text{days}$).
  \item \textbf{Frequency (F):} número total de pedidos.
  \item \textbf{Monetary (M):} gasto total acumulado en euros.
\end{itemize}

Cada dimensión se convierte a una puntuación 1--4 mediante cuartiles calculados con \texttt{numpy.quantile} sobre la población completa de clientes. La recencia se puntúa de forma inversa (menos días $=$ mejor puntuación).

\begin{mypython}[caption={Asignación a segmentos RFM.},label={code:rfm}]
if   sR >= 3 and sF == 4 and sM == 4: segment = "VIP"
elif sF >= 3 and sM >= 3:             segment = "En crecimiento"
elif sR == 1 and sF <= 2:             segment = "En riesgo"
else:                                 segment = "Ocasional"
\end{mypython}

La \textbf{complejidad} del algoritmo es $O(n)$ donde $n$ es el número de clientes, ya que el cálculo de cuartiles y la asignación de puntuaciones son operaciones lineales.

\subsection{Análisis de cesta (Market Basket Analysis)}

El análisis de cesta \cite{agrawal1994} descubre pares de productos comprados frecuentemente juntos. Para cada par $(A, B)$ que co-ocurre en al menos \texttt{min\_support} pedidos, se calculan tres métricas:

\begin{align}
  \text{Soporte}(A, B) &= \frac{|A \cap B|}{N} \label{eq:soporte} \\
  \text{Confianza}(A \to B) &= \frac{|A \cap B|}{|A|} \label{eq:confianza} \\
  \text{Lift}(A \to B) &= \frac{\text{Confianza}(A \to B)}{P(B)} \label{eq:lift}
\end{align}

donde $N$ es el total de pedidos, $|A \cap B|$ el número de pedidos que contienen ambos productos, $|A|$ los pedidos que contienen $A$, y $P(B) = |B|/N$. Un lift $> 1$ indica asociación positiva.

La implementación genera los pares con \texttt{itertools.combinations} y cuenta co-ocurrencias en un solo recorrido. La \textbf{complejidad} es $O(m \cdot k^2)$ donde $m$ es el número de pedidos y $k$ el número medio de productos por pedido.

\subsection{Predicción con suavizado exponencial de Holt}

La predicción de ingresos mensuales utiliza el \textbf{suavizado exponencial doble} de Holt \cite{holt1957}, equivalente a un ARIMA(0,1,1) con tendencia. El modelo mantiene un componente de \textbf{nivel} $L_t$ y un componente de \textbf{tendencia} $T_t$:

\begin{align}
  L_t &= \alpha \cdot y_t + (1 - \alpha)(L_{t-1} + T_{t-1}) \\
  T_t &= \beta (L_t - L_{t-1}) + (1 - \beta) T_{t-1}
\end{align}

con parámetros $\alpha = 0.3$ y $\beta = 0.1$. El pronóstico a $h$ meses vista es:

\begin{equation}
  \hat{y}_{t+h} = L_t + h \cdot T_t
\end{equation}

Los intervalos de predicción al 80\,\% se calculan como $\pm\, 1.28 \cdot \hat{\sigma} \cdot \sqrt{h}$, donde $\hat{\sigma}$ es el RMSE dentro de la muestra. La implementación es en Python puro (librería \texttt{math}), sin dependencias externas.


\section{Integración con LLM (Groq / Llama-3)}

FurniGest incorpora un asistente de inteligencia artificial basado en \textbf{Llama-3.1-8b-instant} \cite{llama3}, ejecutado en la plataforma de inferencia Groq \cite{groq}. La arquitectura se basa en el mecanismo de \textbf{atención} (\emph{attention}) de los modelos Transformer \cite{vaswani2017}: el modelo pre-entrenado se reutiliza para razonamiento analítico mediante \emph{prompt engineering} (\textbf{transfer learning} sin \emph{fine-tuning}).

\textbf{Flujo de una consulta:}
\begin{enumerate}[nosep]
  \item El frontend envía la pregunta del usuario y el rango de fechas a \texttt{POST /api/ai/ask}.
  \item El backend calcula las métricas de negocio (ventas, RFM, cesta, predicción) y las compacta en un JSON reducido ($\leq 120$ puntos de datos).
  \item Se construye un mensaje con un \emph{system prompt} minimalista (idioma, tono, formato JSON para gráficos) y el contexto de datos.
  \item Se llama a la API de Groq (endpoint compatible con OpenAI) con \texttt{temperature=0.2} y \texttt{max\_tokens=1024}.
  \item El backend extrae especificaciones de gráfico del texto de respuesta mediante una expresión regular y las devuelve al frontend.
  \item Si la API no está disponible, se devuelve un resumen determinista calculado directamente sobre las métricas.
\end{enumerate}

La tabla~\ref{tab:llm_params} recoge los parámetros del modelo.

\begin{table}[H]
  \centering
  \caption{Parámetros de configuración del LLM.}
  \label{tab:llm_params}
  \begin{tabular}{lll}
    \toprule
    \textbf{Parámetro} & \textbf{Valor} & \textbf{Justificación} \\
    \midrule
    Modelo       & llama-3.1-8b-instant & 8B parámetros suficientes para razonamiento sobre JSON. \\
    Temperature  & 0.2                  & Salida consistente y factual; evita alucinaciones numéricas. \\
    Timeout      & 60\,s                & Tolerancia para picos de carga en Groq. \\
    \bottomrule
  \end{tabular}
\end{table}


\section{Frontend --- React SPA}

\subsection{Enrutamiento y layout}

React Router 7 \cite{reactrouter} define las rutas de la aplicación. El componente \texttt{App.jsx} actúa como \emph{shell}: renderiza la barra lateral (\texttt{Sidebar.jsx}) una sola vez y conmuta el contenido principal mediante \texttt{<Outlet>}. Las doce rutas del sistema (tabla~\ref{tab:rutas}) se registran con \texttt{createBrowserRouter}.

\begin{table}[H]
  \centering
  \caption{Rutas del frontend.}
  \label{tab:rutas}
  \begin{tabular}{llp{6.5cm}}
    \toprule
    \textbf{Ruta} & \textbf{Componente} & \textbf{Función} \\
    \midrule
    \texttt{/}               & Dashboard        & KPIs, gráfico de ingresos, distribución de estados \\
    \texttt{/ventas/nueva}   & NuevaVenta       & Asistente de creación de pedido paso a paso \\
    \texttt{/clientes}       & ClientesPage     & Tabla de clientes con búsqueda y panel de historial \\
    \texttt{/albaranes}      & AlbaranesPage    & Lista de albaranes con filtro y control de estado \\
    \texttt{/productos}      & ProductosPage    & Productos y proveedores (vista dividida) \\
    \texttt{/movimientos}    & MovimientosPage  & Libro de movimientos con totales mensuales \\
    \texttt{/tendencias}     & Tendencias       & Analíticas, chat IA, exportación PDF \\
    \texttt{/transporte}     & TransportePage   & Gestión de rutas y camiones \\
    \texttt{/banco}          & BancoPage        & Panel de pagos Stripe \\
    \texttt{/incidencias}    & IncidenciasPage  & Gestión de incidencias de entrega \\
    \texttt{/perfil}         & PerfilPage       & Cambio de contraseña y nombre de usuario \\
    \texttt{/personalizacion}& PersonalizacionPage & Tema, paleta, logo, idioma, resumen semanal \\
    \bottomrule
  \end{tabular}
\end{table}

\subsection{Gestión de estado}

El estado global se gestiona con la \textbf{Context API} de React \cite{react}:

\begin{itemize}[nosep]
  \item \texttt{ThemeContext}: modo oscuro/claro y paleta activa (6 opciones), persistidos en \texttt{localStorage}.
  \item \texttt{ConfigContext}: nombre de tienda y logotipo, obtenidos de \texttt{GET /api/config} al montar la aplicación.
\end{itemize}

\subsection{Personalización y accesibilidad}

FurniGest ofrece seis paletas de color (Warm, Slate, Forest, Rose, Ocean, Lavender) y un \emph{toggle} de modo oscuro. La internacionalización se implementa con \textbf{i18next} y \texttt{react-i18next}: la interfaz completa está traducida a español e inglés, con detección del idioma almacenado en \texttt{localStorage}.

En materia de accesibilidad \cite{iso9241}, la interfaz cumple WCAG 2.1 nivel AA:
\begin{itemize}[nosep]
  \item \texttt{aria-current="page"} en el enlace activo de la barra lateral.
  \item \texttt{aria-hidden} en iconos decorativos.
  \item \texttt{htmlFor}/\texttt{id} en formularios.
  \item \texttt{role="switch"} con \texttt{aria-checked} en el \emph{toggle} de modo oscuro.
  \item Contraste mínimo 4.5:1 en todas las paletas (TailwindCSS 4 \cite{tailwindcss}).
\end{itemize}

\subsection{Gráficos y visualización}

Los gráficos se renderizan con \textbf{Chart.js 4} \cite{chartjs} a través de \texttt{react-chartjs-2}. Se utilizan tres tipos de gráfico: \texttt{<Line>} (evolución de ingresos y predicción), \texttt{<Bar>} (top productos, comparativa de períodos) y \texttt{<Pie>} (distribución de estados de albarán). El asistente IA genera dinámicamente especificaciones de gráfico en JSON que el frontend renderiza como componentes Chart.js adicionales.


% #############################################################################
%
%                CAPÍTULO 5: PRUEBAS, CALIDAD Y DESPLIEGUE
%
% #############################################################################
\chapter{Pruebas, calidad y despliegue}
\label{chap:pruebas}

Este capítulo es el núcleo de la ingeniería del software del proyecto. Describe la estrategia de pruebas, el análisis estático, la integración continua, la contenerización y el despliegue en producción.

\section{Estrategia de pruebas}

La suite de pruebas de FurniGest sigue la \textbf{pirámide de pruebas} de Cohn \cite{cohn}: muchas pruebas unitarias en la base, un número moderado de pruebas de integración en el centro, y pocas pruebas end-to-end (E2E) en la cima.

\begin{figure}[H]
  \centering
  \begin{tikzpicture}
    \filldraw[fill=green!20, draw=green!50!black]
      (0,0) -- (6,0) -- (4.5,2) -- (1.5,2) -- cycle;
    \filldraw[fill=yellow!20, draw=yellow!50!black]
      (1.5,2) -- (4.5,2) -- (3.75,3.5) -- (2.25,3.5) -- cycle;
    \filldraw[fill=red!20, draw=red!50!black]
      (2.25,3.5) -- (3.75,3.5) -- (3,5) -- cycle;

    \node at (3,0.9) {\textbf{Unitarias: 439}};
    \node at (3,2.6) {\textbf{Integración}};
    \node at (3,4.0) {\textbf{E2E: 107}};

    \node[right] at (6.3,0.8) {\small pytest + Vitest};
    \node[right] at (6.3,2.5) {\small httpx + TestClient};
    \node[right] at (6.3,4.0) {\small Selenium};
  \end{tikzpicture}
  \caption{Pirámide de pruebas de FurniGest.}
  \label{fig:piramide_tests}
\end{figure}

La tabla~\ref{tab:cobertura_global} resume la cobertura por nivel.

\begin{table}[H]
  \centering
  \caption{Resumen de cobertura por nivel de prueba.}
  \label{tab:cobertura_global}
  \begin{tabular}{llcc}
    \toprule
    \textbf{Nivel} & \textbf{Runner} & \textbf{Tests} & \textbf{Cobertura} \\
    \midrule
    Unitario backend    & pytest 9       & 210  & $>$ 80\,\% \\
    Unitario frontend   & Vitest 4       & 229  & $>$ 80\,\% \\
    End-to-end          & Selenium       & 107  & Flujos completos \\
    \midrule
    \textbf{Total}      &                & \textbf{546} & \\
    \bottomrule
  \end{tabular}
\end{table}


\section{Pruebas unitarias del backend}

\textbf{Stack:} pytest 9 \cite{pytest} + httpx (cliente ASGI para \texttt{TestClient}) + pytest-asyncio + pytest-mock.

\textbf{Aislamiento:} las pruebas utilizan una base de datos SQLite en memoria (\texttt{:memory:}) con \texttt{StaticPool} (conexión única compartida). Todas las tablas se crean y eliminan en cada test mediante un \emph{fixture} \texttt{autouse}. \texttt{PRAGMA foreign\_keys=ON} se aplica mediante un \emph{event listener} de SQLAlchemy para que SQLite reproduzca el comportamiento de PostgreSQL respecto a claves foráneas.

\textbf{Test doubles:} siguiendo la taxonomía de Meszaros \cite{meszaros}, el proyecto utiliza los cinco tipos clásicos de \emph{test double}:

\begin{table}[H]
  \centering
  \caption{Tipos de \emph{test double} utilizados.}
  \label{tab:testdoubles}
  \begin{tabularx}{\textwidth}{llX}
    \toprule
    \textbf{Tipo} & \textbf{Ejemplo} & \textbf{Uso en FurniGest} \\
    \midrule
    Fake  & SQLite in-memory   & Sustituye PostgreSQL para aislar pruebas de la red. \\
    Stub  & Respuesta fija de Groq & \texttt{mocker.patch("...groq\_chat", return\_value=GROQ\_STUB)} \\
    Spy   & Stripe Session.create & \texttt{spy\_create.assert\_called\_once()} \\
    Mock  & \texttt{send\_email\_with\_pdf} & Parche a no-op para evitar conexión SMTP. \\
    Dummy & Email return value & Retorno vacío; el valor no se inspecciona. \\
    \bottomrule
  \end{tabularx}
\end{table}

\begin{mypython}[caption={Configuración de base de datos para pruebas (conftest.py).},label={code:conftest}]
# Motor SQLite en memoria (StaticPool = misma conexion siempre)
engine = create_engine(
    "sqlite:///:memory:",
    connect_args={"check_same_thread": False},
    poolclass=StaticPool,
)

@event.listens_for(engine, "connect")
def _enable_foreign_keys(dbapi_connection, _):
    dbapi_connection.execute("PRAGMA foreign_keys=ON")

@pytest.fixture(autouse=True)
def setup_db():
    """Crea tablas antes de cada test y las borra al terminar."""
    Base.metadata.create_all(bind=engine)
    yield
    Base.metadata.drop_all(bind=engine)
\end{mypython}

La tabla~\ref{tab:tests_backend} detalla los ficheros de test del backend.

\begin{table}[H]
  \centering
  \caption{Tests unitarios del backend.}
  \label{tab:tests_backend}
  \begin{tabular}{lcp{6cm}}
    \toprule
    \textbf{Fichero} & \textbf{Tests} & \textbf{Área cubierta} \\
    \midrule
    test\_clientes.py      & 12 & CRUD + upsert por DNI/email \\
    test\_proveedores.py   &  9 & CRUD completo \\
    test\_productos.py     & 11 & CRUD + búsqueda + violación FK 409 \\
    test\_movimientos.py   & 12 & CRUD + orden por fecha + tipo inválido 422 \\
    test\_albaranes.py     & 26 & Crear, listar, fianza auto, estados, movimiento pendiente, PDF, líneas \\
    test\_analytics.py     & 27 & summary, compare, RFM, cesta, Holt, export/pdf con Groq mockeado \\
    test\_transportes.py   & 14 & Almacén, rutas, asignar, liquidar \\
    test\_stripe.py        & 10 & Checkout, confirm, list (Stripe mockeado) \\
    test\_bank.py          &  3 & Debug info, status, link \\
    test\_auth.py          & 16 & Login OK/fail, me GET/PUT, token expirado/manipulado \\
    test\_configuracion.py & 12 & GET defaults, GET/PUT round-trip, key desconocida \\
    test\_emailer.py       &  7 & html\_to\_text, SMTP path, Resend path \\
    test\_resumen\_semanal.py & 39 & \_get/\_set, \_build\_html, \_run condiciones, Groq error, Holt forecast \\
    test\_incidencias.py   & 12 & Crear, listar, eliminar, estado automático \\
    \midrule
    \textbf{Total}         & \textbf{210} & \\
    \bottomrule
  \end{tabular}
\end{table}


\section{Pruebas unitarias del frontend}

\textbf{Stack:} Vitest 4 \cite{vitest} + \texttt{@testing-library/react} + happy-dom.

\textbf{Aislamiento:}
\begin{itemize}[nosep]
  \item \texttt{globalThis.fetch} se reemplaza por un \texttt{vi.fn()} que devuelve respuestas vacías.
  \item \texttt{ResizeObserver} se \emph{stubbea} (no disponible en happy-dom) para evitar errores de Chart.js.
  \item \texttt{react-chartjs-2} se sustituye por componentes \texttt{<canvas>} ligeros.
  \item Todas las llamadas a \texttt{render()} se envuelven en \texttt{await act(async () => \{ \ldots \})} para que los \texttt{useEffect} se resuelvan antes de las aserciones.
\end{itemize}

La tabla~\ref{tab:tests_frontend} detalla los ficheros de test del frontend.

\begin{table}[H]
  \centering
  \caption{Tests unitarios del frontend.}
  \label{tab:tests_frontend}
  \begin{tabular}{lcp{6cm}}
    \toprule
    \textbf{Fichero} & \textbf{Tests} & \textbf{Área cubierta} \\
    \midrule
    LoginPage.test.jsx          & 18 & Formulario, visibilidad password, login OK/fail, redirect \\
    auth.test.js                & 13 & saveToken, getToken, removeToken, isTokenValid, login \\
    fetchInterceptor.test.js    &  9 & Cabecera auth, URLs externas, 401 $\to$ logout \\
    PersonalizacionPage.test.jsx& 48 & Dark mode, paletas, email, formularios, errores API \\
    PerfilPage.test.jsx         & 17 & Decode JWT, campos password, submit, validación \\
    ThemeContext.test.jsx        & 11 & Valores por defecto, localStorage, setIsDark, setPalette \\
    Sidebar.test.jsx            &  9 & Enlaces nav, clases CSS activas, logout \\
    App.test.jsx                &  5 & Renderizado root, layout, aside + main \\
    Resto (11 ficheros)          & 99 & Dashboard, Albaranes, Banco, Clientes, Movimientos, NuevaVenta, Productos, Tendencias, Transporte, Incidencias, i18n \\
    \midrule
    \textbf{Total}              & \textbf{229} & \\
    \bottomrule
  \end{tabular}
\end{table}


\section{Pruebas end-to-end}

Las pruebas E2E utilizan \textbf{Selenium WebDriver} \cite{selenium} con Chrome en modo \emph{headless}. Las fixtures de pytest proporcionan un navegador con sesión iniciada (\texttt{logged\_in\_browser}, ámbito de sesión) y otro limpio (\texttt{fresh\_driver}, ámbito de módulo). Las pruebas se ejecutan contra la aplicación real (backend en \texttt{localhost:8000}, frontend en \texttt{localhost:5173}).

La tabla~\ref{tab:tests_e2e} lista los 13 ficheros de tests E2E.

\begin{table}[H]
  \centering
  \caption{Tests end-to-end (Selenium).}
  \label{tab:tests_e2e}
  \begin{tabular}{lp{8cm}}
    \toprule
    \textbf{Fichero} & \textbf{Flujos cubiertos} \\
    \midrule
    test\_e2e\_login.py          & Login OK, login fallido, redirect a /login sin token \\
    test\_e2e\_dashboard.py      & Carga de KPIs, gráfico de ingresos, gráfico de estados \\
    test\_e2e\_clientes.py       & Crear, buscar, editar, eliminar cliente \\
    test\_e2e\_albaranes.py      & Crear albarán, cambiar estado, filtrar \\
    test\_e2e\_nueva\_venta.py   & Flujo completo: cliente $\to$ productos $\to$ confirmar \\
    test\_e2e\_productos.py      & CRUD productos y proveedores \\
    test\_e2e\_movimientos.py    & Listar movimientos, crear manual, filtrar por tipo \\
    test\_e2e\_tendencias.py     & Gráficos, chat IA, exportar PDF \\
    test\_e2e\_transporte.py     & Asignar a camión, liquidar ruta \\
    test\_e2e\_banco.py          & Estado Stripe, historial de pagos \\
    test\_e2e\_incidencias.py    & Crear y eliminar incidencia \\
    test\_e2e\_personalizacion.py& Dark mode, paleta, idioma \\
    test\_e2e\_perfil.py         & Cambiar contraseña \\
    \bottomrule
  \end{tabular}
\end{table}


\section{Análisis estático de código}

\subsection{Ruff (Python)}

\textbf{Ruff} \cite{ruff} es un linter y formateador para Python que reemplaza a Flake8, isort y Black con un solo binario. Se aplica sobre todo el directorio \texttt{backend/}:

\begin{itemize}[nosep]
  \item \texttt{ruff check}: verifica cumplimiento PEP~8, importaciones no usadas (F401), sentencias múltiples (E702).
  \item \texttt{ruff format --check}: verifica que el formateo coincide con el estilo normalizado (modo \emph{read-only} en CI).
\end{itemize}

\subsection{ESLint (JavaScript / React)}

ESLint se configura en \texttt{eslint.config.js} (formato plano, ESLint 9) con reglas para React y excepciones para los \emph{globals} de Vitest (\texttt{vi}, \texttt{describe}, \texttt{it}, \texttt{expect}, etc.) y para bloques \texttt{catch \{\}} vacíos (\texttt{allowEmptyCatch: true}).

\subsection{SonarCloud}

FurniGest integra \textbf{SonarCloud} \cite{sonarcloud} como \emph{quality gate}. El análisis se ejecuta en el job \texttt{sonarcloud} del pipeline \texttt{ci.yml}, recibiendo los informes de cobertura generados por pytest (\texttt{coverage-backend.xml}) y Vitest (\texttt{lcov.info}). SonarCloud evalúa:

\begin{itemize}[nosep]
  \item \textbf{Cobertura en código nuevo} $\geq$ 80\,\%.
  \item Bugs, vulnerabilidades y \emph{code smells}.
  \item Si el \emph{quality gate} falla, el despliegue se bloquea.
\end{itemize}

Tanto Ruff, ESLint como SonarCloud actúan como herramientas de \textbf{análisis estático} (\emph{shift-left}): los defectos se detectan antes de que el código llegue a producción, alineándose con la filosofía de desarrollo seguro de OWASP \cite{owasp2021}.


\section{Integración continua (CI)}

El proyecto utiliza dos \textbf{ficheros de workflow de GitHub Actions} \cite{githubactions} que, en conjunto, definen nueve \emph{jobs}:

\begin{table}[H]
  \centering
  \caption{Workflows de GitHub Actions.}
  \label{tab:workflows}
  \begin{tabularx}{\textwidth}{llX}
    \toprule
    \textbf{Fichero} & \textbf{Trigger} & \textbf{Jobs} \\
    \midrule
    \texttt{ci.yml} & PR a \texttt{main} / manual & 7 jobs secuenciales: lint $\to$ backend-tests \& frontend-tests (paralelo) $\to$ SonarCloud $\to$ E2E $\to$ Docker build \& push $\to$ deploy (Vercel + Railway) \\
    \texttt{lint.yml} & push a cualquier rama / manual & 2 jobs independientes: Ruff (Python) y ESLint (JavaScript) \\
    \bottomrule
  \end{tabularx}
\end{table}

El fichero principal (\texttt{ci.yml}) implementa el pipeline completo de CI/CD con siete jobs encadenados:

\begin{figure}[H]
  \centering
  \begin{tikzpicture}[
    pstep/.style={rectangle, draw, rounded corners, minimum width=2.2cm, minimum height=0.8cm, align=center, font=\small},
    >=Stealth
  ]
    \node[pstep, fill=gray!15]    (lint)    {Lint\\(Ruff + ESLint)};
    \node[pstep, fill=green!15, right=0.8cm of lint, yshift=0.6cm]   (tfe)     {Frontend\\tests};
    \node[pstep, fill=green!15, right=0.8cm of lint, yshift=-0.6cm]  (tbe)     {Backend\\tests};
    \node[pstep, fill=yellow!15, right=0.8cm of tfe, yshift=-0.6cm]  (sonar)   {SonarCloud\\gate};
    \node[pstep, fill=orange!15, right=0.8cm of sonar]  (e2e)     {E2E\\Selenium};
    \node[pstep, fill=purple!15, right=0.8cm of e2e]    (docker)  {Docker\\build};
    \node[pstep, fill=blue!15, right=0.8cm of docker]   (deploy)  {Deploy\\Railway+Vercel};

    \draw[->] (lint) |- (tfe);
    \draw[->] (lint) |- (tbe);
    \draw[->] (tfe) -| (sonar);
    \draw[->] (tbe) -| (sonar);
    \draw[->] (sonar) -- (e2e);
    \draw[->] (e2e) -- (docker);
    \draw[->] (docker) -- (deploy);
  \end{tikzpicture}
  \caption{Pipeline CI/CD del workflow \texttt{ci.yml}.}
  \label{fig:pipeline}
\end{figure}

\textbf{Ficheros de configuración \emph{dummy} en CI:} los cinco ficheros Python con secretos están excluidos del repositorio (\texttt{.gitignore}). El paso <<Create dummy config files>> del \emph{workflow} genera módulos con cadenas vacías para desbloquear la cadena de importación de FastAPI. Los valores reales nunca se usan en pruebas ---todas las llamadas externas (SMTP, Stripe, Groq) están mockeadas.


\section{Contenerización con Docker}

Docker Compose \cite{docker} orquesta tres contenedores que se comunican a través de una red bridge aislada.

\begin{figure}[H]
  \centering
  \begin{tikzpicture}[
    container/.style={rectangle, draw, rounded corners, minimum width=3cm, minimum height=1.8cm, align=center, font=\small},
    >=Stealth
  ]
    \node[container, fill=blue!10]   (fe)  {frontend\\{\footnotesize nginx:alpine}\\{\footnotesize :80}};
    \node[container, fill=green!10, right=2cm of fe]  (be)  {backend\\{\footnotesize python:3.12-slim}\\{\footnotesize :8000}};
    \node[container, fill=orange!10, right=2cm of be] (db)  {db\\{\footnotesize postgres:16-alpine}\\{\footnotesize :5432}};

    \draw[<->, thick] (fe) -- node[above, font=\scriptsize] {proxy\_pass /api} (be);
    \draw[<->, thick] (be) -- node[above, font=\scriptsize] {DATABASE\_URL} (db);

    \node[below=0.7cm of db, font=\scriptsize\itshape] {Volume: postgres\_data};
  \end{tikzpicture}
  \caption{Arquitectura de contenedores Docker.}
  \label{fig:docker}
\end{figure}

\begin{table}[H]
  \centering
  \caption{Servicios Docker Compose.}
  \label{tab:docker}
  \begin{tabular}{llcl}
    \toprule
    \textbf{Servicio} & \textbf{Imagen base} & \textbf{Puerto} & \textbf{Función} \\
    \midrule
    db       & postgres:16-alpine & 5432 & Base de datos con healthcheck (\texttt{pg\_isready}) \\
    backend  & python:3.12-slim   & 8000 & API REST (Uvicorn \cite{uvicorn}) \\
    frontend & nginx:alpine       & 80   & SPA + reverse proxy \texttt{/api} \\
    \bottomrule
  \end{tabular}
\end{table}

\textbf{Nginx} actúa como \textbf{proxy inverso}: las peticiones a \texttt{/api/*} se reenvían al contenedor backend; el resto de rutas sirven \texttt{index.html} (\texttt{try\_files \$uri /index.html}), lo que permite que React Router gestione la navegación del lado cliente. Los ficheros de configuración con secretos se montan como volúmenes \emph{read-only} desde el host, evitando que los secretos se incluyan en la imagen Docker.

\section{Despliegue en producción}

\subsection{Infraestructura}

\begin{table}[H]
  \centering
  \caption{Infraestructura de producción.}
  \label{tab:produccion}
  \begin{tabular}{lll}
    \toprule
    \textbf{Componente} & \textbf{Proveedor} & \textbf{Detalle} \\
    \midrule
    Backend + BD  & Railway \cite{railway}  & Contenedor Python + PostgreSQL, región europe-west4 \\
    Frontend      & Vercel \cite{vercel}    & Build Vite estático, CDN edge global \\
    \bottomrule
  \end{tabular}
\end{table}

\subsection{Despliegue continuo (CD)}

Cada Pull Request a \texttt{main} que supera el pipeline completo desencadena dos despliegues automáticos:
\begin{itemize}[nosep]
  \item \textbf{Railway:} el CLI ejecuta \texttt{railway up --detach} tras pasar los tests y el \emph{quality gate}. El script \texttt{start.sh} ejecuta \texttt{alembic upgrade head} antes de iniciar Uvicorn, garantizando que el esquema está sincronizado.
  \item \textbf{Vercel:} el CLI ejecuta \texttt{vercel --prod}, generando un build estático distribuido en CDN.
\end{itemize}

\subsection{Seguridad en producción}

La comunicación en producción está protegida por \textbf{TLS} (Transport Layer Security) \cite{kurose}: Railway aprovisiona y termina el certificado X.509 de forma automática. TLS cifra el tráfico HTTP con \textbf{AES} (cifrado simétrico de sesión), garantizando confidencialidad e integridad de los datos en tránsito.

\textbf{Email en producción:} Railway bloquea las conexiones SMTP salientes (puertos 587/465). El módulo \texttt{emailer.py} detecta la presencia de \texttt{RESEND\_API\_KEY} y conmuta automáticamente a la API HTTP de Resend \cite{resend} sin necesitar cambios en el código ni reinicio.


\section{Migraciones de base de datos (Alembic)}

\textbf{Alembic} \cite{alembic} gestiona la evolución del esquema en producción. A diferencia de \texttt{create\_all()} (que solo crea tablas inexistentes), Alembic genera ficheros de migración versionados con funciones \texttt{upgrade()} y \texttt{downgrade()}, registrando el estado actual en una tabla \texttt{alembic\_version} dentro de PostgreSQL.

\begin{mypython}[caption={start.sh --- arranque en producción con migración automática.},label={code:startsh}]
#!/bin/bash
alembic upgrade head
exec uvicorn backend.app.main:app \
    --host 0.0.0.0 --port "${PORT:-8000}"
\end{mypython}


% #############################################################################
%
%                        CAPÍTULO 6: RESULTADOS
%
% #############################################################################
\chapter{Resultados}
\label{chap:resultados}

\section{Evaluación heurística de usabilidad}

La interfaz de FurniGest se ha evaluado utilizando las \textbf{diez heurísticas de Nielsen} \cite{nielsen}. La tabla~\ref{tab:nielsen} recoge cada heurística y su implementación en el sistema.

\begin{longtable}{p{4.5cm}p{8.5cm}}
  \caption{Evaluación heurística de Nielsen aplicada a FurniGest.} \label{tab:nielsen} \\
  \toprule
  \textbf{Heurística} & \textbf{Implementación} \\
  \midrule
  \endfirsthead
  \toprule
  \textbf{Heurística} & \textbf{Implementación} \\
  \midrule
  \endhead
  Visibilidad del estado del sistema & Notificaciones \emph{toast} (sileo) confirman cada operación; los estados del albarán se muestran con badges de color diferenciado (ámbar, azul, violeta, verde, rojo). \\
  Coincidencia entre el sistema y el mundo real & La terminología refleja el vocabulario del negocio: albarán, fianza, ruta, camión, liquidación. \\
  Control y libertad del usuario & Las acciones destructivas (eliminar cliente, eliminar incidencia) requieren confirmación explícita; filtros y búsquedas se pueden limpiar con un clic. \\
  Consistencia y estándares & TailwindCSS 4 aplica una paleta de colores, tipografía y espaciado uniforme en todas las pantallas. \\
  Prevención de errores & Los formularios validan campos en el cliente antes del envío; Pydantic rechaza peticiones malformadas en el servidor con código 422. \\
  Reconocimiento antes que recuerdo & Los iconos y etiquetas del menú lateral son siempre visibles; las tablas incluyen cabeceras descriptivas. \\
  Flexibilidad y eficiencia de uso & Atajos de teclado implícitos (\emph{Enter} para buscar), barra de búsqueda rápida en listados. \\
  Diseño estético y minimalista & Cada pantalla muestra solo las columnas y acciones relevantes; sin información superflua. \\
  Ayudar a los usuarios a reconocer, diagnosticar y recuperarse de errores & Los mensajes de error son descriptivos (<<La contraseña actual no es correcta>>, <<Cliente con ese DNI ya existe>>). \\
  Ayuda y documentación & La documentación interactiva de la API está accesible en \texttt{/docs} (Swagger UI); la interfaz incluye textos de ayuda en formularios. \\
  \bottomrule
\end{longtable}


\section{Métricas de calidad del software}

La tabla~\ref{tab:metricas} recoge las métricas de calidad reportadas por SonarCloud y la suite de pruebas.

\begin{table}[H]
  \centering
  \caption{Métricas de calidad del software.}
  \label{tab:metricas}
  \begin{tabular}{lc}
    \toprule
    \textbf{Métrica} & \textbf{Valor} \\
    \midrule
    Quality Gate SonarCloud     & Passed \\
    Cobertura backend (pytest)  & $>$ 80\,\% \\
    Cobertura frontend (Vitest) & $>$ 80\,\% \\
    Bugs detectados             & 0 \\
    Vulnerabilidades            & 0 \\
    Tests unitarios backend     & 210 \\
    Tests unitarios frontend    & 229 \\
    Tests E2E (Selenium)        & 107 \\
    Tests totales               & 546 \\
    Endpoints REST              & 61 \\
    Entidades de dominio        & 11 \\
    Ficheros de workflow CI/CD  & 2 (9 jobs) \\
    \bottomrule
  \end{tabular}
\end{table}


\section{Capturas de pantalla}

% NOTA: Las capturas deben tomarse del sistema en ejecución y añadirse
% como ficheros de imagen (.png o .pdf) en el directorio memoria/.
% A continuación se reserva el espacio con la estructura adecuada.

\textit{Las siguientes capturas de pantalla deben añadirse desde el sistema en ejecución:}

\begin{enumerate}[nosep]
  \item \textbf{Login:} formulario de inicio de sesión con logo de tienda.
  \item \textbf{Dashboard:} tarjetas de KPIs, gráfico de ingresos (línea) y distribución de estados (tarta).
  \item \textbf{Nueva Venta:} asistente paso a paso con búsqueda de cliente y adición de productos.
  \item \textbf{Clientes:} tabla con búsqueda y panel lateral de historial de pedidos.
  \item \textbf{Albaranes:} lista filtrable con control de estado inline.
  \item \textbf{Productos:} vista dividida con productos a la izquierda y proveedores a la derecha.
  \item \textbf{Movimientos:} libro de ingresos y egresos con totales mensuales.
  \item \textbf{Tendencias:} gráficos de analíticas, comparativa de períodos, chat IA.
  \item \textbf{Transporte:} columna de almacén y columnas por camión.
  \item \textbf{Banco:} panel de pagos Stripe con historial de cobros.
  \item \textbf{Personalización:} selector de paleta, modo oscuro, idioma.
  \item \textbf{SonarCloud:} dashboard del quality gate.
\end{enumerate}

% Ejemplo de cómo incluir una captura cuando esté disponible:
% \begin{figure}[H]
%   \centering
%   \includegraphics[width=0.85\textwidth]{captura_dashboard.png}
%   \caption{Pantalla principal (Dashboard) de FurniGest.}
%   \label{fig:cap_dashboard}
% \end{figure}


\section{URLs de producción}

La tabla~\ref{tab:urls} recoge las URLs del sistema desplegado.

\begin{table}[H]
  \centering
  \caption{URLs del sistema en producción.}
  \label{tab:urls}
  \begin{tabularx}{\textwidth}{lX}
    \toprule
    \textbf{Servicio} & \textbf{URL} \\
    \midrule
    Frontend      & \url{https://tfg-five-drab.vercel.app} \\
    Backend API   & \url{https://supportive-nurturing-production-66dc.up.railway.app} \\
    Swagger (docs)& \url{https://supportive-nurturing-production-66dc.up.railway.app/docs} \\
    \bottomrule
  \end{tabularx}
\end{table}


% #############################################################################
%
%              CAPÍTULO 7: CONCLUSIONES Y TRABAJOS FUTUROS
%
% #############################################################################
\chapter{Conclusiones y trabajos futuros}

\section{Conclusiones}

Este Trabajo Fin de Grado ha cumplido el objetivo de diseñar, implementar y desplegar en producción un sistema de información de gestión para una tienda de muebles. A continuación se evalúa el grado de cumplimiento de cada objetivo específico:

\begin{description}[nosep]
  \item[OBJ-01 (CRUD):] Cumplido. Clientes, productos y proveedores disponen de operaciones CRUD completas con búsqueda y filtrado.
  \item[OBJ-02 (Albarán):] Cumplido. La máquina de cinco estados funciona con transiciones controladas y efectos secundarios automáticos.
  \item[OBJ-03 (Movimientos):] Cumplido. Los movimientos se generan automáticamente al crear albaranes, al entregar pedidos y al liquidar rutas.
  \item[OBJ-04 (Logística):] Cumplido. El módulo de transporte permite asignar albaranes a camiones, visualizar manifiestos y liquidar rutas con factura PDF.
  \item[OBJ-05 (Analíticas):] Cumplido. El motor implementa segmentación RFM, análisis de cesta y predicción de Holt, todo accesible desde la pantalla de tendencias.
  \item[OBJ-06 (IA):] Cumplido. El asistente basado en Llama-3 responde preguntas de negocio y genera gráficos a partir de datos en tiempo real.
  \item[OBJ-07 (Stripe):] Cumplido. El flujo de pago con Stripe Checkout funciona con registro automático de ingresos e idempotencia.
  \item[OBJ-08 (Pruebas):] Cumplido. Más de 540 pruebas automatizadas cubren los tres niveles de la pirámide de pruebas.
  \item[OBJ-09 (CI/CD):] Cumplido. Dos ficheros de workflow de GitHub Actions (con nueve jobs en total) garantizan que cada cambio pasa por pruebas, linting, quality gate y despliegue.
  \item[OBJ-10 (Producción):] Cumplido. El sistema está accesible públicamente en Railway (backend) y Vercel (frontend).
\end{description}

El enfoque orientado a \textbf{producción real} ha sido determinante: trabajar con restricciones reales (Railway bloquea SMTP, Vercel requiere builds estáticos, SonarCloud exige cobertura $\geq$ 80\,\%) obligó a tomar decisiones de ingeniería que no habrían surgido en un prototipo académico.

La inversión en \textbf{ingeniería del software} ---pruebas, CI/CD, análisis estático, quality gate--- ha permitido mantener la base de código mantenible y confiable a medida que crecían las funcionalidades. Cada Pull Request se ha fusionado con la garantía de que pasa el pipeline completo.

\section{Trabajos futuros}

La tabla~\ref{tab:trabajos_futuros} propone líneas de evolución del sistema.

\begin{table}[H]
  \centering
  \caption{Trabajos futuros propuestos.}
  \label{tab:trabajos_futuros}
  \begin{tabularx}{\textwidth}{lX}
    \toprule
    \textbf{Trabajo futuro} & \textbf{Descripción} \\
    \midrule
    Roles y permisos granulares   & Diferenciar permisos de vendedor, repartidor y administrador con listas de control de acceso. \\
    Gestión de stock e inventario & Registrar entradas y salidas de almacén en tiempo real, con alertas de stock mínimo. \\
    App móvil (PWA o React Native)& Acceso optimizado para repartidores en ruta con capacidades offline. \\
    Caché Redis para analíticas   & Reducir la carga de consultas frecuentes al motor de analíticas. \\
    WebSockets para notificaciones& Actualizar el dashboard en tiempo real cuando se crea un pedido o cambia un estado. \\
    Módulo de presupuestos        & Generar y enviar presupuestos previos al pedido formal. \\
    Soporte multi-tienda          & Gestionar varias sucursales desde una única instancia. \\
    Pruebas de carga (Locust)     & Validar el rendimiento del sistema bajo concurrencia elevada. \\
    Conciliación bancaria automática & Madurar la integración con CaixaBank PSD2 para detectar automáticamente cobros recibidos y conciliarlos con albaranes pendientes de pago. \\
    \bottomrule
  \end{tabularx}
\end{table}


% #############################################################################
%                            BIBLIOGRAFÍA
% #############################################################################

\phantomsection
\addcontentsline{toc}{chapter}{Bibliografía}

\footnotesize{
\bibliographystyle{IEEEtran}
\bibliography{bibliografia}
}

\raggedbottom
\afterpage{\blankpage}
\newpage


% #############################################################################
%                            APÉNDICES
% #############################################################################

\appendix

\phantomsection
\addcontentsline{toc}{chapter}{Apéndices}

\mbox{}
\vfill
\begin{center}
\begin{Huge}
\textbf{Apéndices}
\end{Huge}
\end{center}
\vfill
\mbox{}
\thispagestyle{empty}

\newpage
\mbox{}
\thispagestyle{empty}
\newpage


% =====================================================
%  APÉNDICE A: MANUAL DE INSTALACIÓN
% =====================================================
\chapter{Manual de instalación y configuración}
\label{apx:instalacion}

\section{Requisitos previos}

\begin{itemize}[nosep]
  \item Python 3.12 con entorno virtual (\texttt{.venv/}).
  \item PostgreSQL 16 con una base de datos \texttt{TFG} creada.
  \item Node.js 20+.
  \item Docker Desktop (opcional, para ejecución con contenedores).
\end{itemize}

\section{Instalación local}

\subsection{Backend}

\begin{verbatim}
# Desde la raíz del repositorio
.venv\Scripts\activate          # Windows (PowerShell)
pip install -r requirements.txt
uvicorn backend.app.main:app --reload
# API disponible en http://localhost:8000
# Documentación en http://localhost:8000/docs
\end{verbatim}

\textbf{Ficheros de configuración:} crear los siguientes ficheros en \texttt{backend/app/} con los valores reales (ver \texttt{.env.example} en la raíz):
\begin{itemize}[nosep]
  \item \texttt{auth\_config.py} --- JWT secret key y configuración.
  \item \texttt{stripe\_config.py} --- claves Stripe.
  \item \texttt{ia\_settings.py} --- API key de Groq.
  \item \texttt{settings\_email.py} --- configuración SMTP / Resend.
  \item \texttt{bank\_settings.py} --- configuración bancaria.
\end{itemize}

\subsection{Frontend}

\begin{verbatim}
cd frontend
npm install
npm run dev
# App disponible en http://localhost:5173
\end{verbatim}

\section{Instalación con Docker Compose}

\begin{verbatim}
docker compose up --build
# Frontend: http://localhost
# API:      http://localhost/api/...
# Backend:  http://localhost:8000
\end{verbatim}

Los ficheros de configuración se montan como volúmenes read-only. Deben existir en el host antes de ejecutar \texttt{docker compose up}.

\section{Datos de demostración}

Al arrancar por primera vez, \texttt{seed.py} genera: 18 proveedores, 99 productos, 100 clientes, $\sim$150 albaranes y $\sim$260 movimientos. El seed es idempotente: si ya hay datos, no se ejecuta.

\textbf{Cuenta por defecto:} usuario \texttt{admin}, contraseña \texttt{admin123}. Cambiar en producción.


% =====================================================
%  APÉNDICE B: CATÁLOGO DE ENDPOINTS REST
% =====================================================
\chapter{Catálogo de endpoints REST}
\label{apx:endpoints}

\begin{longtable}{llp{6cm}}
  \caption{Catálogo completo de endpoints.} \label{tab:endpoints_full} \\
  \toprule
  \textbf{Método} & \textbf{Ruta} & \textbf{Descripción} \\
  \midrule
  \endfirsthead
  \toprule
  \textbf{Método} & \textbf{Ruta} & \textbf{Descripción} \\
  \midrule
  \endhead
  \multicolumn{3}{l}{\textbf{Clientes}} \\
  GET    & /api/clientes/get     & Listar todos \\
  GET    & /api/clientes/get/\{id\} & Obtener uno \\
  POST   & /api/clientes/post    & Crear \\
  PUT    & /api/clientes/put/\{id\} & Actualizar \\
  DELETE & /api/clientes/delete/\{id\} & Eliminar \\
  \midrule
  \multicolumn{3}{l}{\textbf{Productos}} \\
  GET    & /api/productos/get    & Listar todos (con proveedor) \\
  GET    & /api/productos/get/\{id\} & Obtener uno \\
  GET    & /api/productos/search & Buscar por nombre \\
  POST   & /api/productos/post   & Crear \\
  PUT    & /api/productos/put/\{id\} & Actualizar \\
  DELETE & /api/productos/delete/\{id\} & Eliminar \\
  \midrule
  \multicolumn{3}{l}{\textbf{Proveedores}} \\
  GET    & /api/proveedores/get  & Listar todos \\
  GET    & /api/proveedores/get/\{id\} & Obtener uno \\
  POST   & /api/proveedores/post & Crear \\
  PUT    & /api/proveedores/put/\{id\} & Actualizar \\
  DELETE & /api/proveedores/delete/\{id\} & Eliminar \\
  \midrule
  \multicolumn{3}{l}{\textbf{Albaranes}} \\
  POST   & /api/albaranes/post   & Crear pedido completo \\
  GET    & /api/albaranes/get    & Listar todos \\
  GET    & /api/albaranes/get/\{id\} & Obtener uno \\
  GET    & /api/albaranes/by-cliente/\{id\} & Pedidos de un cliente \\
  PUT    & /api/albaranes/put/\{id\} & Actualizar campos editables \\
  PUT    & /api/albaranes/\{id\}/items & Editar líneas del albarán \\
  PATCH  & /api/albaranes/\{id\}/estado & Avanzar estado \\
  GET    & /api/albaranes/\{id\}/pdf & Descargar albarán en PDF \\
  \midrule
  \multicolumn{3}{l}{\textbf{Movimientos}} \\
  GET    & /api/movimientos/get  & Listar (filtro por tipo) \\
  GET    & /api/movimientos/get/\{id\} & Obtener uno \\
  POST   & /api/movimientos/post & Crear manualmente \\
  PUT    & /api/movimientos/put/\{id\} & Actualizar \\
  DELETE & /api/movimientos/delete/\{id\} & Eliminar \\
  \midrule
  \multicolumn{3}{l}{\textbf{Transporte}} \\
  GET    & /api/transporte/almacen & Albaranes en ALMACEN \\
  GET    & /api/transporte/ruta  & Albaranes en RUTA \\
  GET    & /api/transporte/rutas & Agrupados por camión \\
  POST   & /api/transporte/ruta/asignar & Asignar a camión \\
  POST   & /api/transporte/ruta/quitar  & Quitar de camión \\
  POST   & /api/transporte/ruta/pendiente & Marcar como pendiente \\
  POST   & /api/transporte/ruta/\{id\}/liquidar & Liquidar ruta (genera factura PDF) \\
  \midrule
  \multicolumn{3}{l}{\textbf{Analíticas}} \\
  GET    & /api/analytics/summary  & Métricas completas \\
  GET    & /api/analytics/compare  & Comparativa de períodos \\
  GET    & /api/analytics/predict  & Predicción Holt \\
  GET    & /api/analytics/export/pdf & PDF de tendencias \\
  \midrule
  \multicolumn{3}{l}{\textbf{IA}} \\
  POST   & /api/ai/ask  & Pregunta puntual \\
  POST   & /api/ai/chat & Conversación multi-turno \\
  \midrule
  \multicolumn{3}{l}{\textbf{Stripe}} \\
  GET    & /api/stripe/status    & Estado de configuración \\
  POST   & /api/stripe/checkout  & Crear sesión de pago \\
  POST   & /api/stripe/confirm   & Confirmar pago \\
  GET    & /api/stripe/checkouts & Historial de cobros \\
  \midrule
  \multicolumn{3}{l}{\textbf{Incidencias}} \\
  GET    & /api/incidencias/get  & Listar todas \\
  GET    & /api/incidencias/get/\{id\} & Obtener una \\
  POST   & /api/incidencias/post & Crear incidencia \\
  DELETE & /api/incidencias/\{id\} & Eliminar (revierte albarán) \\
  \midrule
  \multicolumn{3}{l}{\textbf{Banca (CaixaBank PSD2)}} \\
  GET    & /api/bank/caixa/\_debug   & Info de depuración \\
  GET    & /api/bank/caixa/status    & Estado de conexión \\
  POST   & /api/bank/caixa/link      & Iniciar enlace con banco \\
  GET    & /api/bank/caixa/callback  & Callback OAuth \\
  GET    & /api/bank/caixa/accounts  & Listar cuentas \\
  \midrule
  \multicolumn{3}{l}{\textbf{Auth}} \\
  POST   & /api/auth/login & Obtener JWT \\
  GET    & /api/auth/me    & Obtener perfil \\
  PUT    & /api/auth/me    & Actualizar perfil \\
  \midrule
  \multicolumn{3}{l}{\textbf{Config}} \\
  GET    & /api/config        & Obtener configuración \\
  PUT    & /api/config/\{key\} & Actualizar clave \\
  \midrule
  \multicolumn{3}{l}{\textbf{Health}} \\
  GET    & /health & Sonda de salud \\
  \bottomrule
\end{longtable}


% Fin del documento
\end{document}
